\documentclass[11pt,a4paper]{article}

% ==================== TEMEL PAKETLER (pdflatex) ====================
\usepackage[T1]{fontenc}
\usepackage[utf8]{inputenc}
\usepackage{lmodern}
\usepackage[a4paper,margin=2cm,top=2.4cm,bottom=2.4cm]{geometry}
\usepackage{xcolor}
\usepackage{colortbl}
\usepackage{array}
\usepackage{booktabs}
\usepackage{longtable}
\usepackage{tabularx}
\usepackage{listings}
\usepackage{fancyhdr}
\usepackage{hyperref}

% ==================== RENK PALETİ ====================
\definecolor{anaMavi}{HTML}{1B3A5B}
\definecolor{vurguMavi}{HTML}{2E6DA4}
\definecolor{acikMavi}{HTML}{EAF2F9}
\definecolor{turuncu}{HTML}{D9701E}
\definecolor{yesil}{HTML}{2E7D32}
\definecolor{acikYesil}{HTML}{E8F5E9}
\definecolor{kirmizi}{HTML}{C0392B}
\definecolor{mor}{HTML}{6C3483}
\definecolor{acikMor}{HTML}{F3E9F7}
\definecolor{acikGri}{HTML}{F4F5F7}
\definecolor{ortaGri}{HTML}{6B7280}
\definecolor{koyuGri}{HTML}{2D2D2D}
\definecolor{kodArka}{HTML}{FBFBFA}

% ==================== HYPERREF ====================
\hypersetup{
  colorlinks=true,
  linkcolor={[HTML]{2E6DA4}},
  urlcolor={[HTML]{2E6DA4}},
  pdftitle={İleri Seviye C ve Linux Sistem Programlama Rehberi},
  pdfauthor={C Rehberi},
  bookmarksnumbered=true
}

% ==================== BAŞLIK STİLLERİ (titlesec olmadan) ====================
\setcounter{secnumdepth}{0}
\newcounter{secc}
\newcounter{subsecc}[secc]
\newcounter{subsubsecc}[subsecc]
\newcommand{\gsec}[1]{\par\phantomsection\stepcounter{secc}%
  \setcounter{subsecc}{0}%
  \vspace{18pt}{\color{anaMavi}\Large\bfseries\thesecc.\hspace{0.5em}#1}\par
  \vspace{2pt}{\color{turuncu}\rule{\linewidth}{1.6pt}}\par\vspace{8pt}%
  \addcontentsline{toc}{section}{\protect\numberline{\thesecc}#1}\markboth{#1}{#1}}
\newcommand{\gsub}[1]{\par\phantomsection\stepcounter{subsecc}%
  \setcounter{subsubsecc}{0}%
  \vspace{11pt}{\color{vurguMavi}\large\bfseries\thesecc.\thesubsecc\hspace{0.5em}#1}\par\vspace{4pt}%
  \addcontentsline{toc}{subsection}{\protect\numberline{\thesecc.\thesubsecc}#1}}
\newcommand{\gsubsub}[1]{\par\phantomsection\stepcounter{subsubsecc}%
  \vspace{7pt}{\color{koyuGri}\normalsize\bfseries #1}\par\vspace{2pt}}

% Kısım (Part) başlığı
\newcommand{\gpart}[2]{\clearpage
  \begin{center}\vspace*{1.5cm}
  {\color{turuncu}\rule{0.7\linewidth}{1.5pt}}\\[10pt]
  {\color{ortaGri}\Large KISIM #1}\\[8pt]
  {\color{anaMavi}\Huge\bfseries #2}\\[10pt]
  {\color{turuncu}\rule{0.7\linewidth}{1.5pt}}
  \end{center}\vspace{1cm}}

% ==================== BAŞLIK / ALTBİLGİ ====================
\setlength{\headheight}{14pt}
\pagestyle{fancy}
\fancyhf{}
\renewcommand{\headrulewidth}{0.4pt}
\renewcommand{\footrulewidth}{0.4pt}
\fancyhead[L]{\small\color{ortaGri}İleri C \& Linux Sistem Programlama}
\fancyhead[R]{\small\color{ortaGri}\leftmark}
\fancyfoot[C]{\small\color{ortaGri}\thepage}

% ==================== KOD LİSTELEME (C) ====================
\lstdefinestyle{cstyle}{
  language=C,
  basicstyle=\ttfamily\small\color{koyuGri},
  keywordstyle=\color{vurguMavi}\bfseries,
  commentstyle=\color{yesil}\itshape,
  stringstyle=\color{turuncu},
  numberstyle=\tiny\color{ortaGri},
  numbers=left,
  numbersep=8pt,
  backgroundcolor=\color{kodArka},
  frame=single,
  framesep=6pt,
  rulecolor=\color{ortaGri!40},
  breaklines=true,
  showstringspaces=false,
  tabsize=4,
  columns=fullflexible,
  keepspaces=true,
  extendedchars=true,
  literate={ı}{{\i}}1{İ}{{\.I}}1{ş}{{\c s}}1{Ş}{{\c S}}1{ğ}{{\u g}}1{Ğ}{{\u G}}1{ç}{{\c c}}1{Ç}{{\c C}}1{ö}{{\"o}}1{Ö}{{\"O}}1{ü}{{\"u}}1{Ü}{{\"U}}1,
  morekeywords={size_t,ssize_t,pid_t,uid_t,gid_t,off_t,mode_t,socklen_t,
    uint8_t,uint16_t,uint32_t,uint64_t,int8_t,int16_t,int32_t,int64_t,
    intptr_t,uintptr_t,ptrdiff_t,NULL,bool,true,false,restrict,inline,
    _Generic,_Atomic,_Noreturn,pthread_t,pthread_mutex_t,pthread_cond_t,
    sem_t,sig_atomic_t,va_list,FILE}
}
\lstset{style=cstyle}

% Kabuk/terminal stili
\lstdefinestyle{shstyle}{
  basicstyle=\ttfamily\small\color{koyuGri},
  backgroundcolor=\color{koyuGri!6},
  frame=single,framesep=6pt,rulecolor=\color{ortaGri!40},
  breaklines=true,showstringspaces=false,tabsize=4,columns=fullflexible,
  keepspaces=true,
  literate={ı}{{\i}}1{İ}{{\.I}}1{ş}{{\c s}}1{Ş}{{\c S}}1{ğ}{{\u g}}1{Ğ}{{\u G}}1{ç}{{\c c}}1{Ç}{{\c C}}1{ö}{{\"o}}1{Ö}{{\"O}}1{ü}{{\"u}}1{Ü}{{\"U}}1,
}

% ==================== ÖZEL KUTULAR (tcolorbox olmadan) ====================
\newcommand{\callout}[4]{\par\smallskip\noindent\begingroup
  \setlength{\fboxsep}{8pt}\setlength{\fboxrule}{0.8pt}%
  \fcolorbox{#1}{#2}{\parbox{\dimexpr\linewidth-2\fboxsep-2\fboxrule\relax}{%
    {\bfseries\color{#1}#3}\par\smallskip #4}}\endgroup\par\smallskip}
\newcommand{\bilgi}[1]{\callout{vurguMavi}{acikMavi}{Bilgi}{#1}}
\newcommand{\ipucu}[1]{\callout{yesil}{acikYesil}{İpucu}{#1}}
\newcommand{\uyari}[1]{\callout{kirmizi}{kirmizi!7}{Dikkat}{#1}}
\newcommand{\ornek}[2][Örnek]{\callout{ortaGri!90!black}{acikGri}{#1}{#2}}
\newcommand{\notk}[2][Not]{\callout{mor}{acikMor}{#1}{#2}}

% Yardımcı komutlar
\newcommand{\code}[1]{\texttt{\color{koyuGri}#1}}
\newcommand{\hd}[1]{\texttt{\color{koyuGri}\textless #1\textgreater}}
\renewcommand{\arraystretch}{1.35}

% ==================== BAŞLANGIÇ ====================
\begin{document}

% ---------- KAPAK ----------
\begin{titlepage}
\centering
\vspace*{1.4cm}
{\color{turuncu}\rule{\textwidth}{2pt}}\\[8pt]
{\color{anaMavi}\Huge\bfseries İleri Seviye C}\\[10pt]
{\color{vurguMavi}\LARGE ve Linux Sistem Programlama Rehberi}\\[8pt]
{\color{ortaGri}\large İşaretçilerden Sistem Çağrılarına Kapsamlı Başvuru}\\[6pt]
{\color{turuncu}\rule{\textwidth}{2pt}}\\[1.4cm]

\begin{center}
\fcolorbox{vurguMavi}{acikMavi}{\parbox{0.86\textwidth}{\centering\large
Bu belge C dilinin ileri konularını (işaretçiler, birleşimler, fonksiyon
işaretçileri, dinamik bellek, önişlemci) ve Linux sistem programlamayı
(dosya G/Ç, süreçler, sinyaller, IPC, iş parçacıkları, soketler)
\textbf{çok sayıda çalışan kod örneğiyle} adım adım anlatır.\vspace{4pt}}}
\end{center}

\vspace{0.8cm}
\begin{center}
\fcolorbox{ortaGri!50}{acikGri}{\parbox{0.86\textwidth}{\small
\textbf{Kapsam --- Kısım 1 (İleri C):} bellek modeli \textbullet\ saklama sınıfları
\textbullet\ işaretçiler (temelden ileriye) \textbullet\ işaretçi aritmetiği
\textbullet\ diziler \& bellek düzeni \textbullet\ dinamik bellek (malloc/free)
\textbullet\ yapılar \& hizalama \textbullet\ birleşimler \& tagged union
\textbullet\ enum \& bit alanları \textbullet\ fonksiyon işaretçileri \& callback
\textbullet\ değişken argümanlı fonksiyonlar \textbullet\ önişlemci \& makrolar
\textbullet\ jenerik programlama \textbullet\ bağlı liste / yığın / kuyruk.\vspace{4pt}

\textbf{Kapsam --- Kısım 2 (Linux Sistem Programlama):} sistem çağrıları \& errno
\textbullet\ düşük seviye dosya G/Ç \textbullet\ dosya sistemi \& stat
\textbullet\ süreçler (fork/exec/wait) \textbullet\ sinyaller \textbullet\ borular
\& FIFO \textbullet\ System V \& POSIX IPC \textbullet\ paylaşımlı bellek
\textbullet\ pthreads \& senkronizasyon \textbullet\ TCP/UDP soketleri
\textbullet\ mmap \textbullet\ zaman \& zamanlayıcılar.\vspace{4pt}}}
\end{center}

\vfill
{\color{ortaGri}\large Hazırlanma Tarihi: 9 Temmuz 2026}\\[4pt]
{\color{ortaGri} GCC / glibc \textbullet\ POSIX \textbullet\ Linux x86-64}
\vspace*{0.6cm}
\end{titlepage}

% ---------- İÇİNDEKİLER ----------
\tableofcontents
\thispagestyle{empty}
\clearpage

% #####################################################################
\gpart{1}{İleri Seviye C Programlama}
% #####################################################################

% =====================================================================
\gsec{Giriş: C'nin Bellek Modeli ve Derleme Süreci}
% =====================================================================

C, donanıma yakın çalışan, belleği doğrudan yöneten bir dildir. İleri C'yi
anlamanın anahtarı, programın çalışırken belleği nasıl kullandığını
zihinde canlandırabilmektir. Bu bölüm, sonraki tüm konuların üzerine
kurulacağı temeli atar.

\gsub{Derleme Aşamaları}

Bir \code{.c} dosyası çalıştırılabilir hale gelene kadar dört aşamadan geçer:

\begin{center}
\begin{tabularx}{\textwidth}{@{}>{\bfseries\color{anaMavi}}l l >{\raggedright\arraybackslash}X@{}}
\toprule
\rowcolor{acikMavi}\color{anaMavi}\textbf{Aşama} & \color{anaMavi}\textbf{Araç} & \color{anaMavi}\textbf{Ne yapar?}\\
\midrule
Önişleme & cpp & \code{\#include}, \code{\#define}, \code{\#ifdef} işlenir; saf C üretilir.\\
\rowcolor{acikGri} Derleme & cc1 & C kodu mimariye özgü assembly'ye çevrilir (\code{.s}).\\
Assembly & as & Assembly, nesne koduna (\code{.o}) çevrilir.\\
\rowcolor{acikGri} Bağlama & ld & Nesne dosyaları + kütüphaneler birleştirilip çalıştırılabilir üretilir.\\
\bottomrule
\end{tabularx}
\end{center}

\begin{lstlisting}[style=shstyle]
$ gcc -E  main.c -o main.i    # yalnızca önişleme
$ gcc -S  main.c -o main.s    # assembly üret
$ gcc -c  main.c -o main.o    # nesne dosyası üret
$ gcc     main.c -o main      # tüm aşamalar + bağlama

# Rehber boyunca önerilen derleme bayrakları:
$ gcc -std=c11 -Wall -Wextra -g -O2 main.c -o main
#   -std=c11 : C11 standardı     -Wall -Wextra : tüm uyarılar
#   -g       : hata ayıklama bilgisi   -O2 : optimizasyon
\end{lstlisting}

\gsub{Bir Programın Bellek Düzeni}

Çalışan bir C programı, sanal bellekte birbirinden ayrı bölgelere yerleşir.
Bu düzeni bilmek işaretçileri, \code{static} değişkenleri ve segfault'ları
anlamanın temelidir.

\begin{center}
\begin{tabular}{|>{\centering\arraybackslash}p{0.34\textwidth}|>{\raggedright\arraybackslash}p{0.56\textwidth}|}
\hline
\rowcolor{anaMavi}\multicolumn{2}{|c|}{\color{white}\textbf{Yüksek Adresler}}\\
\hline
\cellcolor{acikMavi}\textbf{Stack (Yığın)} & Yerel değişkenler, fonksiyon parametreleri, dönüş adresleri. Aşağı doğru büyür.\\
\hline
\multicolumn{2}{|c|}{\cellcolor{acikGri}\itshape $\downarrow$ \quad boş alan \quad $\uparrow$}\\
\hline
\cellcolor{acikYesil}\textbf{Heap (Öbek)} & \code{malloc}/\code{calloc} ile ayrılan dinamik bellek. Yukarı doğru büyür.\\
\hline
\cellcolor{acikMor}\textbf{BSS} & İlklendirilmemiş global/static değişkenler (sıfırlanır).\\
\hline
\cellcolor{acikMor}\textbf{Data} & İlklendirilmiş global/static değişkenler.\\
\hline
\cellcolor{turuncu!20}\textbf{Text (Kod)} & Makine kodu; salt-okunur.\\
\hline
\rowcolor{anaMavi}\multicolumn{2}{|c|}{\color{white}\textbf{Düşük Adresler}}\\
\hline
\end{tabular}
\end{center}

\ornek[Değişkenlerin nerede oturduğunu görme]{Aşağıdaki program her değişken
türünün adresini yazdırır. Adresleri karşılaştırdığında segmentlerin ayrıldığını
görürsün.}
\begin{lstlisting}
#include <stdio.h>
#include <stdlib.h>

int global_ilk = 5;      // Data segmenti
int global_sifir;        // BSS segmenti

int main(void) {
    int yerel = 1;                     // Stack
    static int statik = 2;             // Data
    int *dinamik = malloc(sizeof(int)); // Heap

    printf("Text (main)   : %p\n", (void*)main);
    printf("Data (global) : %p\n", (void*)&global_ilk);
    printf("BSS  (sifir)  : %p\n", (void*)&global_sifir);
    printf("Heap (malloc) : %p\n", (void*)dinamik);
    printf("Stack (yerel) : %p\n", (void*)&yerel);
    printf("static        : %p\n", (void*)&statik);

    free(dinamik);
    return 0;
}
\end{lstlisting}

\bilgi{Stack otomatik yönetilir (fonksiyon bitince temizlenir), heap ise
\emph{senin sorumluluğundadır}: \code{malloc} ile aldığını \code{free} ile
geri vermelisin. Bu ayrım, ileri C'nin en kritik konusudur.}

% =====================================================================
\gsec{Bellek, Saklama Sınıfları ve Nitelikler}
% =====================================================================

\gsub{Saklama Sınıfları (Storage Classes)}

Saklama sınıfı, bir değişkenin \emph{ömrünü} (lifetime) ve \emph{görünürlüğünü}
(scope) belirler.

\begin{center}
\begin{tabularx}{\textwidth}{@{}>{\ttfamily\color{anaMavi}}l l l >{\raggedright\arraybackslash}X@{}}
\toprule
\rowcolor{acikMavi}\color{anaMavi}\textbf{Anahtar} & \color{anaMavi}\textbf{Ömür} & \color{anaMavi}\textbf{Görünürlük} & \color{anaMavi}\textbf{Açıklama}\\
\midrule
auto & blok & blok & Varsayılan; yerel değişkenler. Neredeyse hiç yazılmaz.\\
\rowcolor{acikGri} register & blok & blok & Register'da tutulması \emph{önerilir} (modern derleyiciler yok sayar).\\
static & program & sınırlı & Değeri çağrılar arası korunur; dosya kapsamında dışarı kapalıdır.\\
\rowcolor{acikGri} extern & program & global & Başka dosyada tanımlı değişkene/fonksiyona bildirim.\\
\bottomrule
\end{tabularx}
\end{center}

\ornek[static ile durum saklama]{Fonksiyon içi \code{static} değişken, çağrılar
arasında değerini korur.}
\begin{lstlisting}
#include <stdio.h>

int sayac(void) {
    static int n = 0;   // yalnızca ilk çağrıda ilklenir
    return ++n;
}

int main(void) {
    printf("%d %d %d\n", sayac(), sayac(), sayac()); // 1 2 3 (bkz. not)
    return 0;
}
\end{lstlisting}

\uyari{Yukarıda \code{printf} argümanlarının değerlendirilme sırası
\emph{tanımsızdır}; farklı derleyiciler "3 2 1" de yazabilir. Yan etkili
ifadeleri tek bir \code{printf}'e sıralı biçimde koymaktan kaçın.}

\gsub{const, volatile ve restrict Nitelikleri}

\begin{itemize}
\item \code{const}: değer değiştirilemez (salt-okunur). Derleyici zorlar.
\item \code{volatile}: değer her erişimde bellekten okunur; derleyici
      optimize edip önbelleğe almaz. Donanım register'ları ve sinyal
      bayrakları için şarttır.
\item \code{restrict} (C99): "bu işaretçi, gösterdiği belleğe erişen tek
      işaretçidir" sözü verir; derleyici agresif optimizasyon yapar.
\end{itemize}

\begin{lstlisting}
const double PI = 3.14159;         // değiştirilemez
volatile sig_atomic_t bayrak = 0;  // sinyal handler'ı değiştirebilir

// restrict: src ve dst çakışmaz sözü -> daha hızlı kod
void kopyala(int *restrict dst, const int *restrict src, size_t n) {
    for (size_t i = 0; i < n; i++) dst[i] = src[i];
}
\end{lstlisting}

\notk[const'un yeri anlamı değiştirir]{\code{const int *p} $\to$ p'nin
gösterdiği \emph{değer} sabittir. \code{int *const p} $\to$ p'nin kendisi
(adres) sabittir. \code{const int *const p} $\to$ ikisi de sabittir.}

% =====================================================================
\gsec{İşaretçiler: Temelden İleri Seviyeye}
% =====================================================================

İşaretçi, bir \emph{bellek adresini} tutan değişkendir. C'nin gücünün ve
tehlikesinin kaynağıdır. Bu bölümü sindirmek, ileri C'nin yarısını halletmek
demektir.

\gsub{Temel Kavram: Adres ve Yönelme}

\begin{itemize}
\item \code{\&x} --- x değişkeninin \emph{adresini} verir (address-of).
\item \code{*p} --- p'nin gösterdiği adresteki \emph{değeri} verir (dereference).
\item \code{int *p} --- "p, bir \code{int}'i gösteren işaretçidir".
\end{itemize}

\begin{lstlisting}
#include <stdio.h>
int main(void) {
    int x = 42;
    int *p = &x;        // p, x'in adresini tutar

    printf("x  = %d\n", x);      // 42
    printf("&x = %p\n", (void*)&x);
    printf("p  = %p\n", (void*)p);   // &x ile aynı
    printf("*p = %d\n", *p);     // 42  (p'nin gösterdiği değer)

    *p = 100;           // x'i işaretçi üzerinden değiştir
    printf("x  = %d\n", x);      // 100
    return 0;
}
\end{lstlisting}

\bilgi{İşaretçinin \emph{tipi} önemlidir: \code{int *} bir int gösterir,
\code{double *} bir double. Tip, \code{*p} ile kaç bayt okunacağını ve
işaretçi aritmetiğinin adım boyunu belirler.}

\gsub{NULL, Wild ve Dangling İşaretçiler}

\begin{itemize}
\item \textbf{NULL işaretçi:} hiçbir yeri göstermez (\code{p = NULL}).
      Yönelmek (\code{*p}) segfault üretir. Her zaman kontrol et.
\item \textbf{Wild (başıboş) işaretçi:} ilklenmemiş; rastgele bir adres tutar.
      Çok tehlikelidir. Bildirir bildirmez ilklendir.
\item \textbf{Dangling (asılı) işaretçi:} serbest bırakılmış (\code{free})
      ya da kapsamı bitmiş belleği gösterir.
\end{itemize}

\begin{lstlisting}
int *wild;              // TEHLIKE: rastgele adres
int *safe = NULL;       // guvenli baslangic

int *dangling = malloc(sizeof(int));
free(dangling);         // bellek geri verildi
// *dangling = 5;       // HATA: dangling pointer -> tanimsiz davranis
dangling = NULL;        // free'den sonra NULL'la -> guvenli
\end{lstlisting}

\uyari{Serbest bıraktıktan sonra işaretçiyi \code{NULL}'a ata. Böylece
tekrar \code{free} (double-free) veya erişim (use-after-free) hataları
segfault'a dönüşür ve fark edilir; sessiz bellek bozulması yaşanmaz.}

\gsub{İşaretçi Aritmetiği}

İşaretçiye tamsayı eklemek, adresi \emph{tip boyutu kadar} kaydırır. Yani
\code{p + 1}, bir sonraki \code{int}'i gösterir --- 4 bayt ileriyi
(int için), 1 bayt değil.

\begin{lstlisting}
#include <stdio.h>
int main(void) {
    int a[5] = {10, 20, 30, 40, 50};
    int *p = a;            // p, a[0]'i gosterir (dizi -> isaretcie cozulur)

    printf("%d\n", *p);       // 10
    printf("%d\n", *(p + 1)); // 20  (p+1 -> a[1], 4 bayt ileri)
    printf("%d\n", *(p + 3)); // 40
    p++;                      // artik a[1]'i gosterir
    printf("%d\n", *p);       // 20

    // iki isaretci farki: aradaki ELEMAN sayisi
    int *q = &a[4];
    printf("fark = %ld\n", (long)(q - p)); // 3
    return 0;
}
\end{lstlisting}

\notk[Dizi indeksleme aslında işaretçi aritmetiğidir]{\code{a[i]} ifadesi
tam olarak \code{*(a + i)} demektir. Bu yüzden ilginç ama geçerli olan
\code{i[a]} yazımı da \code{a[i]} ile aynıdır! Çünkü toplama değişmelidir.}

\gsub{Diziler ve İşaretçiler: Aynı Değiller}

Diziler çoğu bağlamda ilk elemanı gösteren işaretçiye "çözülür" (decay),
ama \emph{özdeş değildir}.

\begin{center}
\begin{tabularx}{\textwidth}{@{}>{\raggedright\arraybackslash}X >{\raggedright\arraybackslash}X@{}}
\toprule
\rowcolor{acikMavi}\color{anaMavi}\textbf{Dizi \texttt{int a[5]}} & \color{anaMavi}\textbf{İşaretçi \texttt{int *p}}\\
\midrule
\code{sizeof(a)} = 20 (tüm dizi) & \code{sizeof(p)} = 8 (yalnız adres)\\
\rowcolor{acikGri}\code{a} yeniden atanamaz & \code{p} yeniden atanabilir\\
\code{\&a} tipi \code{int(*)[5]} & \code{\&p} tipi \code{int**}\\
\bottomrule
\end{tabularx}
\end{center}

\begin{lstlisting}
#include <stdio.h>
void fonksiyon(int dizi[]) {          // aslinda int *dizi
    // BURADA sizeof(dizi) == sizeof(int*) == 8, dizi boyutu DEGIL!
    printf("fonksiyon icinde: %zu\n", sizeof(dizi)); // 8
}
int main(void) {
    int a[5];
    printf("main icinde: %zu\n", sizeof(a));  // 20
    fonksiyon(a);
    return 0;
}
\end{lstlisting}

\uyari{Bir dizi fonksiyona geçince boyut bilgisi kaybolur (işaretçiye çözülür).
Bu yüzden dizi boyutunu \emph{ayrı bir parametre} olarak geçmelisin:
\code{void f(int *dizi, size\_t n)}.}

\gsub{İşaretçiden İşaretçiye (Double Pointer)}

Bir işaretçinin adresini tutan işaretçidir: \code{int **pp}. Kullanım alanları:
bir fonksiyonun çağıranın işaretçisini değiştirmesi ve işaretçi dizileri
(örn. \code{char *argv[]}).

\begin{lstlisting}
#include <stdio.h>
#include <stdlib.h>

// Bir fonksiyonun cagiranin isaretcisini degistirmesi icin
// isaretcinin ADRESI (int**) gerekir.
void tahsis_et(int **pp, int deger) {
    *pp = malloc(sizeof(int));  // cagiranin isaretcisine yaz
    **pp = deger;               // o bellege deger yaz
}

int main(void) {
    int *p = NULL;
    tahsis_et(&p, 99);          // p'nin ADRESINI gecir
    printf("%d\n", *p);         // 99
    free(p);
    return 0;
}
\end{lstlisting}

\ornek[İşaretçi dizisi: string tablosu]{\code{char *} dizisi, farklı uzunlukta
stringleri verimli tutar.}
\begin{lstlisting}
#include <stdio.h>
int main(void) {
    const char *gunler[] = {"Pzt", "Sal", "Car", "Per", "Cum"};
    size_t n = sizeof(gunler) / sizeof(gunler[0]);
    for (size_t i = 0; i < n; i++)
        printf("%zu -> %s\n", i, gunler[i]);
    return 0;
}
\end{lstlisting}

\gsub{void İşaretçileri (Genel İşaretçi)}

\code{void *} herhangi bir tipin adresini tutabilir; ama yönelmeden önce
somut bir tipe \emph{cast} edilmelidir. \code{malloc} bu yüzden \code{void *}
döndürür.

\begin{lstlisting}
#include <stdio.h>
void yazdir(void *veri, char tip) {
    switch (tip) {
        case 'i': printf("%d\n",  *(int*)veri);    break;
        case 'd': printf("%g\n",  *(double*)veri); break;
        case 'c': printf("%c\n",  *(char*)veri);   break;
    }
}
int main(void) {
    int i = 5; double d = 3.14; char c = 'A';
    yazdir(&i, 'i');
    yazdir(&d, 'd');
    yazdir(&c, 'c');
    return 0;
}
\end{lstlisting}

\uyari{\code{void *} üzerinde işaretçi aritmetiği standart C'de \emph{yasaktır}
(boyutu bilinmez). GCC bir eklenti olarak 1 bayt kabul eder, ama taşınabilir
kod için önce \code{char *}'a cast et.}

\gsub{const ve İşaretçiler --- Dört Kombinasyon}

\begin{lstlisting}
int x = 10, y = 20;
int *p1 = &x;              // her sey degisebilir
const int *p2 = &x;        // *p2 sabit (deger), p2 degisebilir
int *const p3 = &x;        // p3 sabit (adres), *p3 degisebilir
const int *const p4 = &x;  // ikisi de sabit

// *p2 = 5;   // HATA: gosterilen deger sabit
p2 = &y;      // OK: isaretci baska yeri gosterebilir
*p3 = 5;      // OK: deger degistirilebilir
// p3 = &y;   // HATA: isaretci sabit
\end{lstlisting}

\ipucu{Okuma kuralı: işaretçiyi \emph{sağdan sola} oku. \code{int *const p}
= "p is a const pointer to int". \code{const int *p} = "p is a pointer to
const int".}

% =====================================================================
\gsec{İleri İşaretçi Teknikleri}
% =====================================================================

Bu bölüm, temel işaretçilerin ötesine geçer: karmaşık bildirimler, diziyi
gösteren işaretçiler, hizalama, takma ad (aliasing) kuralları, kernel tarzı
gömülü veri yapıları, opak işaretçiler ve bağlam taşıyan callback'ler. Bunlar
gerçek sistem kodunun (Linux çekirdeği, glibc, gömülü sistemler) günlük
araçlarıdır.

\gsub{Karmaşık Bildirimleri Okuma: Sağ-Sol Kuralı}

C bildirimlerini çözmek için \emph{isimden başla, önce sağa sonra sola} git;
parantezler önceliği değiştirir. \code{()} ve \code{[]} operatörleri \code{*}'dan
önce gelir.

\begin{center}
\begin{tabularx}{\textwidth}{@{}>{\ttfamily\color{anaMavi}}l >{\raggedright\arraybackslash}X@{}}
\toprule
\rowcolor{acikMavi}\color{anaMavi}\textbf{Bildirim} & \color{anaMavi}\textbf{Okunuşu}\\
\midrule
int *p[10] & 10 elemanlı, her biri \code{int*} olan \emph{dizi}.\\
\rowcolor{acikGri} int (*p)[10] & 10 elemanlı bir \code{int} dizisini gösteren \emph{işaretçi}.\\
int *f(void) & \code{int*} döndüren fonksiyon.\\
\rowcolor{acikGri} int (*f)(void) & \code{int} döndüren fonksiyonu gösteren işaretçi.\\
int (*fp[4])(int) & Her biri fonksiyon işaretçisi olan 4'lük dizi (jump table).\\
\rowcolor{acikGri} int (*(*p)[5])(void) & 5'lik fonksiyon-işaretçisi dizisini gösteren işaretçi.\\
char **argv & \code{char*} dizisinin ilk elemanını gösteren işaretçi.\\
\bottomrule
\end{tabularx}
\end{center}

\ipucu{\code{cdecl} aracı bu bildirimleri Türkçe/İngilizce açıklar:
\code{echo 'explain int (*p)[10]' | cdecl}. \code{typedef} kullanmak, karmaşık
tipleri okunur parçalara böler ve en iyi çözümdür.}

\gsub{Diziyi Gösteren İşaretçi: \texttt{int (*p)[N]}}

Diziyi gösteren işaretçi, iki boyutlu dizilerde satır satır ilerlemenin doğru
yoludur. \code{p + 1}, \emph{tüm bir satır} kadar (N eleman) ilerler.

\begin{lstlisting}
#include <stdio.h>
int main(void) {
    int matris[3][4] = {
        { 1,  2,  3,  4},
        { 5,  6,  7,  8},
        { 9, 10, 11, 12},
    };

    // 'satir', 4 int'lik bir diziyi gosterir. p++ -> sonraki SATIR.
    int (*satir)[4] = matris;          // matris -> &matris[0]
    for (int i = 0; i < 3; i++) {
        for (int j = 0; j < 4; j++)
            printf("%3d ", (*satir)[j]);  // (*satir) = matris[i]
        printf("\n");
        satir++;                       // 4 int (16 bayt) ileri: sonraki satir
    }
    return 0;
}
\end{lstlisting}

\uyari{\code{int (*p)[4]} ile \code{int *p[4]} tamamen farklıdır. İlki bir
işaretçidir (8 bayt), ikincisi 4 işaretçilik bir dizidir (32 bayt). Parantez
şart.}

\gsub{Çok Boyutlu Diziler ve İşaretçiye Çözülme}

\code{int a[3][4]} bellekte \emph{bitişik} 12 int'tir (satır-öncelikli). Bir
fonksiyona geçince ilk boyut düşer ve \code{int (*)[4]}'e çözülür --- ikinci
boyut korunmalıdır.

\begin{lstlisting}
#include <stdio.h>
// 2B diziyi dogru gecirme: satir sayisi ayri, sutun tipte gomulu
void yazdir(int satir, int sutun, int m[][sutun]) {  // VLA parametresi (C99)
    for (int i = 0; i < satir; i++) {
        for (int j = 0; j < sutun; j++)
            printf("%d ", m[i][j]);      // m[i][j] = *(*(m+i)+j)
        printf("\n");
    }
}
int main(void) {
    int a[2][3] = {{1,2,3},{4,5,6}};
    yazdir(2, 3, a);                     // int(*)[3] olarak gecer
    return 0;
}
\end{lstlisting}

\notk[İndekslemenin açılımı]{\code{a[i][j]} ifadesi
\code{*(*(a + i) + j)} demektir: önce i. satırın adresi, sonra o satırda j.
eleman. Derleyici satır-öncelikli düzende \code{i * sutun + j} ofsetini hesaplar.}

\gsub{İşaretçi Takma Adı (Aliasing) ve Strict Aliasing Kuralı}

Derleyici, \emph{farklı tipteki} iki işaretçinin aynı belleği göstermediğini
varsayarak optimizasyon yapar (strict aliasing). Bunu ihlal etmek sessiz
hatalara yol açar.

\begin{lstlisting}
#include <stdint.h>
#include <string.h>

float ornek(void) {
    float f = 1.0f;

    // YANLIS: strict aliasing ihlali -> tanimsiz davranis
    // uint32_t u = *(uint32_t *)&f;

    // DOGRU: memcpy ile tip donusumu (derleyici bunu optimize eder)
    uint32_t u;
    memcpy(&u, &f, sizeof u);        // bitleri guvenle kopyala
    u |= 0x80000000;                 // isaret bitini set et (negatifle)
    memcpy(&f, &u, sizeof f);
    return f;                        // -1.0f
}
\end{lstlisting}

\bilgi{\code{char *}, \code{unsigned char *} ve \code{void *} her tipe takma ad
olabilir --- kural onlar için geçerli değildir. Bu yüzden \code{memcpy}
(bayt bazlı) her zaman güvenlidir. Tip cezalandırmanın (type punning) diğer
geçerli yolu, daha önce gördüğümüz \code{union}'dır.}

\gsub{Bellek Hizalama: alignof, alignas, aligned\_alloc}

Her tipin bir hizalama gereksinimi vardır (bir \code{int} genelde 4'ün katı
adreste durur). C11 bunu açıkça kontrol etmeyi sağlar --- SIMD, DMA ve
donanım tamponları için gereklidir.

\begin{lstlisting}
#include <stdio.h>
#include <stdlib.h>
#include <stdalign.h>

int main(void) {
    printf("int hizasi   : %zu\n", alignof(int));      // genelde 4
    printf("double hizasi: %zu\n", alignof(double));   // genelde 8

    // 64 baytlik hizali degisken (ornegin cache satiri)
    alignas(64) char tampon[128];
    printf("tampon adresi %% 64 = %ld\n", (long)((uintptr_t)tampon % 64)); // 0

    // Hizali dinamik ayirma (boyut, hizalamanin kati olmali)
    void *p = aligned_alloc(64, 256);  // 64-hizali, 256 bayt
    if (p) {
        printf("aligned_alloc %% 64 = %ld\n", (long)((uintptr_t)p % 64)); // 0
        free(p);
    }
    return 0;
}
// Derleme: gcc -std=c11 bu.c
\end{lstlisting}

\gsub{uintptr\_t: İşaretçi--Tamsayı Dönüşümü ve Etiketleme}

\code{<stdint.h>}'deki \code{uintptr\_t}, bir işaretçiyi güvenle tutabilen
tamsayı tipidir. Hizalanmış işaretçilerin kullanılmayan alt bitlerine bayrak
gömmek (pointer tagging) ileri bir tekniktir.

\begin{lstlisting}
#include <stdio.h>
#include <stdint.h>
int main(void) {
    int veri = 42;
    int *p = &veri;

    // 4-hizali isaretcinin alt 2 biti daima 0'dir -> etiket icin kullan
    uintptr_t etiketli = (uintptr_t)p | 1u;    // alt bite bayrak koy

    int  bayrak = etiketli & 1u;               // bayragi oku
    int *temiz  = (int *)(etiketli & ~(uintptr_t)3u); // isaretciyi geri al

    printf("bayrak=%d, deger=%d\n", bayrak, *temiz);  // bayrak=1, deger=42
    return 0;
}
\end{lstlisting}

\uyari{Etiketleme yalnızca hizalamanın garanti ettiği bitlerde yapılabilir ve
işaretçiyi kullanmadan önce bitler mutlaka temizlenmelidir. Yanlış kullanım
segfault üretir; bu teknik yalnızca performans-kritik kodda tercih edilir.}

\gsub{offsetof ve container\_of: Gömülü (Intrusive) Veri Yapıları}

Linux çekirdeğinin klasik hilesi: bir yapı üyesinin adresinden, onu içeren
\emph{dış} yapının adresini bulmak. Bu, tek bir liste kodunu tüm tiplerde
kullanmayı sağlar.

\begin{lstlisting}
#include <stdio.h>
#include <stddef.h>          // offsetof

// Bir uye isaretcisinden, kapsayan yapiya ulas:
#define container_of(ptr, type, member) \
    ((type *)((char *)(ptr) - offsetof(type, member)))

struct liste_dugum { struct liste_dugum *sonraki; };  // genel liste dugumu

struct gorev {              // gercek veri: liste dugumunu ICINE gomer
    int id;
    struct liste_dugum dugum;
};

int main(void) {
    struct gorev g = { .id = 7 };
    struct liste_dugum *n = &g.dugum;        // yalnizca dugumu tutuyoruz

    // Dugumden geri kapsayan goreve ulas:
    struct gorev *geri = container_of(n, struct gorev, dugum);
    printf("id = %d\n", geri->id);           // 7
    printf("offset(dugum) = %zu\n", offsetof(struct gorev, dugum));
    return 0;
}
\end{lstlisting}

\bilgi{\code{offsetof(tip, uye)}, bir üyenin yapının başından kaç bayt sonra
başladığını verir. \code{container\_of} bu ofseti işaretçiden çıkararak dış
yapıya ulaşır. Bu desen "intrusive" liste/ağaç yapılarının ve
\code{list\_head} API'sinin temelidir.}

\gsub{Opak İşaretçiler: Bilgi Gizleme ve Soyut Veri Tipleri}

Bir yapının \emph{tanımını} başlıkta gizleyip yalnızca işaretçisini (handle)
sunmak, C'de kapsülleme (encapsulation) sağlar. Kullanıcı içeriği göremez,
yalnızca API fonksiyonlarını çağırır.

\begin{lstlisting}
// ====== yigin.h (kullaniciya acik) ======
typedef struct Yigin Yigin;      // OPAK: tanim gizli, sadece isim
Yigin *yigin_olustur(void);
void   yigin_it(Yigin *, int);
int    yigin_cek(Yigin *, int *);
void   yigin_yok_et(Yigin *);

// ====== yigin.c (uygulama, gizli) ======
#include "yigin.h"
#include <stdlib.h>
struct Yigin {                   // tam tanim yalnizca BURADA gorunur
    int   *veri;
    size_t adet, kapasite;
};
Yigin *yigin_olustur(void) {
    Yigin *y = calloc(1, sizeof *y);   // kullanici sizeof(Yigin) bilemez
    return y;
}
// ... yigin_it / yigin_cek / yigin_yok_et ...
\end{lstlisting}

\ipucu{Opak işaretçi, ikili uyumluluğu (ABI) korur: yapının iç düzenini
değiştirsen bile kullanıcı kodunu yeniden derlemek gerekmez. \code{FILE},
\code{DIR}, \code{pthread\_t} gibi standart tipler bu yüzden opaktır.}

\gsub{İleri Fonksiyon İşaretçileri: Bağlam Taşıyan Callback}

Gerçek callback'ler genellikle bir \code{void *userdata} parametresiyle
"bağlam" (context) taşır --- C'de kapanış (closure) benzetiminin yoludur.

\begin{lstlisting}
#include <stdio.h>

// Callback: her eleman + kullanici verisi alir
typedef void (*Islem)(int eleman, void *userdata);

void her_biri(const int *dizi, size_t n, Islem fn, void *userdata) {
    for (size_t i = 0; i < n; i++)
        fn(dizi[i], userdata);
}

// userdata'yi bir toplam biriktirici olarak kullan
void biriktir(int e, void *userdata) {
    *(long *)userdata += e;
}

int main(void) {
    int a[] = {1, 2, 3, 4, 5};
    long toplam = 0;
    her_biri(a, 5, biriktir, &toplam);   // baglam: &toplam
    printf("toplam = %ld\n", toplam);    // 15
    return 0;
}
\end{lstlisting}

\gsubsub{Fonksiyon İşaretçisi Döndüren Fonksiyon}

\begin{lstlisting}
#include <stdio.h>
typedef int (*Islem)(int, int);
int topla(int a, int b) { return a + b; }
int carp (int a, int b) { return a * b; }

// typedef ile temiz: bir Islem dondurur
Islem secim(char op) { return (op == '+') ? topla : carp; }

// typedef'siz ham sozdizimi (ayni sey):
// int (*secim(char op))(int, int) { ... }

int main(void) {
    Islem f = secim('*');
    printf("%d\n", f(6, 7));   // 42
    return 0;
}
\end{lstlisting}

\gsub{restrict ile Optimizasyon ve Örtüşme (Overlap)}

\code{restrict}, işaretçilerin çakışmadığı sözünü verince derleyici gereksiz
yeniden yüklemeleri (reload) atlar. Söz yalanlanırsa (bloklar örtüşürse)
davranış tanımsızdır --- \code{memcpy} bunu varsayar, \code{memmove} varsaymaz.

\begin{lstlisting}
#include <stddef.h>
// restrict YOK: derleyici a[i] her dongude b'nin degistirmis olabilecegini
// varsayar ve tekrar okur -> yavas.
void olcekle_yavas(int *a, const int *b, size_t n, int k) {
    for (size_t i = 0; i < n; i++) a[i] = b[i] * k;
}
// restrict VAR: a ve b cakismaz sozu -> vektorlestirme mumkun.
void olcekle_hizli(int *restrict a, const int *restrict b, size_t n, int k) {
    for (size_t i = 0; i < n; i++) a[i] = b[i] * k;
}
\end{lstlisting}

\gsub{Çok Katmanlı const ve \texttt{const char **} Tuzağı}

\begin{lstlisting}
// Sezgiye aykiri: char** -> const char** donusumu GECERSIZDIR.
// Cunku izin verilseydi const'lugu delmek mumkun olurdu.
void f(const char **p);
char *dizi[1];
// f(dizi);          // UYARI/HATA: char** -> const char** guvenli degil

// Dogru imza: hem eleman hem isaretci const olmali
void g(const char *const *p);   // bu char** kabul eder
\end{lstlisting}

\ipucu{Çok katmanlı işaretçilerde \code{const}, yalnızca \emph{en dış}
seviyede otomatik güvenli eklenir. İç seviyelerde de const isteniyorsa
imzayı buna göre yaz: \code{const char *const *}.}

% =====================================================================
\gsec{Dinamik Bellek Yönetimi}
% =====================================================================

Derleme anında boyutu bilinmeyen veriler için heap'ten çalışma zamanında
bellek ayrılır. Bu, C'nin en güçlü ama en hataya açık alanıdır.

\gsub{malloc, calloc, realloc, free}

\begin{center}
\begin{tabularx}{\textwidth}{@{}>{\ttfamily\color{anaMavi}}l >{\raggedright\arraybackslash}X@{}}
\toprule
\rowcolor{acikMavi}\color{anaMavi}\textbf{Fonksiyon} & \color{anaMavi}\textbf{Görev}\\
\midrule
malloc(n) & n bayt ayırır; içerik \emph{çöp} (ilklenmemiş). Başarısızsa NULL.\\
\rowcolor{acikGri} calloc(k,n) & k$\times$n bayt ayırır ve \emph{sıfırlar}. Taşma kontrolü yapar.\\
realloc(p,n) & p'yi n bayta büyütür/küçültür; içeriği korur, taşıyabilir.\\
\rowcolor{acikGri} free(p) & p ile ayrılan belleği geri verir. \code{free(NULL)} güvenlidir.\\
\bottomrule
\end{tabularx}
\end{center}

\begin{lstlisting}
#include <stdio.h>
#include <stdlib.h>

int main(void) {
    size_t n = 5;
    int *dizi = malloc(n * sizeof(int));   // ayir
    if (dizi == NULL) {                     // HER ZAMAN kontrol et
        perror("malloc");
        return 1;
    }
    for (size_t i = 0; i < n; i++) dizi[i] = (int)(i * i);

    // Buyut: 5 -> 10 eleman
    int *yeni = realloc(dizi, 10 * sizeof(int));
    if (yeni == NULL) {                     // realloc basarisizsa
        free(dizi);                         // eski blok hala gecerli
        return 1;
    }
    dizi = yeni;                            // dogru kullanim
    for (size_t i = 5; i < 10; i++) dizi[i] = (int)(i * i);

    for (size_t i = 0; i < 10; i++) printf("%d ", dizi[i]);
    printf("\n");

    free(dizi);                             // geri ver
    dizi = NULL;                            // dangling'i onle
    return 0;
}
\end{lstlisting}

\uyari{\code{p = realloc(p, n)} yaygın ama tehlikeli bir kalıptır: realloc
NULL dönerse eski işaretçiyi kaybedip \emph{bellek sızıntısı} yaratırsın.
Yukarıdaki gibi geçici bir değişken kullan.}

\gsub{Yaygın Bellek Hataları}

\begin{center}
\begin{tabularx}{\textwidth}{@{}>{\bfseries\color{kirmizi}}l >{\raggedright\arraybackslash}X@{}}
\toprule
\rowcolor{kirmizi!12}\color{kirmizi}\textbf{Hata} & \color{kirmizi}\textbf{Açıklama}\\
\midrule
Memory leak & \code{free} unutulur; program bellek tüketmeye devam eder.\\
\rowcolor{acikGri} Double free & Aynı blok iki kez \code{free} edilir; heap bozulur.\\
Use-after-free & Serbest bırakılmış belleğe erişim.\\
\rowcolor{acikGri} Buffer overflow & Ayrılan sınırın dışına yazma.\\
Uninitialized read & \code{malloc}'tan gelen çöp değerin okunması.\\
\bottomrule
\end{tabularx}
\end{center}

\ipucu{Bu hataları \textbf{Valgrind} ve \textbf{AddressSanitizer} yakalar:
\newline\code{valgrind --leak-check=full ./program}
\newline\code{gcc -fsanitize=address -g program.c \&\& ./a.out}}

\ornek[2B dizi (matris) dinamik ayırma]{Satır sayısı çalışma zamanında
belliyse, işaretçi dizisi + her satır için ayrı blok kullanılır.}
\begin{lstlisting}
#include <stdlib.h>
int **matris_olustur(size_t satir, size_t sutun) {
    int **m = malloc(satir * sizeof(int *));   // satir isaretcisi
    if (!m) return NULL;
    for (size_t i = 0; i < satir; i++) {
        m[i] = malloc(sutun * sizeof(int));    // her satirin verisi
        if (!m[i]) {                            // hata: onceki satirlari temizle
            for (size_t j = 0; j < i; j++) free(m[j]);
            free(m);
            return NULL;
        }
    }
    return m;
}
void matris_yok_et(int **m, size_t satir) {
    for (size_t i = 0; i < satir; i++) free(m[i]); // once satirlar
    free(m);                                        // sonra isaretci dizisi
}
\end{lstlisting}

% =====================================================================
\gsec{Yapılar (struct), Hizalama ve typedef}
% =====================================================================

\code{struct}, farklı tipteki verileri tek bir birim altında toplar. İleri
kullanımda hizalama (alignment) ve dolgu (padding) kritiktir.

\gsub{Tanım, Erişim ve typedef}

\begin{lstlisting}
#include <stdio.h>
#include <string.h>

struct Nokta { int x, y; };          // etiketli yapi

typedef struct {                      // isimsiz + typedef
    char ad[32];
    int  yas;
    double maas;
} Calisan;

int main(void) {
    struct Nokta n = {3, 4};
    printf("%d,%d\n", n.x, n.y);      // nokta erisimi

    Calisan c;
    strcpy(c.ad, "Ada");
    c.yas = 30; c.maas = 5000.0;

    Calisan *pc = &c;
    printf("%s %d\n", pc->ad, pc->yas); // isaretci uzerinden -> operatoru
    return 0;
}
\end{lstlisting}

\notk[Nokta vs Ok operatörü]{\code{yapi.alan} $\to$ doğrudan erişim.
\code{isaretci->alan} $\to$ işaretçi üzerinden erişim; \code{(*isaretci).alan}
ile birebir aynıdır.}

\gsub{Bellek Hizalama ve Dolgu (Padding)}

Derleyici, her üyeyi kendi boyutunun katı olan bir adrese yerleştirir
(hizalama). Bu yüzden yapının boyutu, üyelerin toplamından \emph{büyük}
olabilir --- aradaki boşluğa dolgu (padding) denir.

\begin{lstlisting}
#include <stdio.h>
struct Kotu {       // kotu siralama -> cok dolgu
    char  a;        // 1 bayt + 3 dolgu
    int   b;        // 4 bayt
    char  c;        // 1 bayt + 3 dolgu
};                  // toplam: 12 bayt

struct Iyi {        // iyi siralama (buyukten kucuge) -> az dolgu
    int   b;        // 4 bayt
    char  a;        // 1 bayt
    char  c;        // 1 bayt + 2 dolgu
};                  // toplam: 8 bayt

int main(void) {
    printf("Kotu: %zu, Iyi: %zu\n", sizeof(struct Kotu), sizeof(struct Iyi));
    return 0;       // Kotu: 12, Iyi: 8
}
\end{lstlisting}

\ipucu{Bellekten tasarruf için yapı üyelerini \emph{büyükten küçüğe} sırala.
Yapıyı sıkıştırmak için \code{\#pragma pack(1)} veya
\code{\_\_attribute\_\_((packed))} kullanılabilir --- ama hizasız erişim
bazı mimarilerde yavaşlar/çöker; ağ protokolleri için dikkatle kullan.}

\gsub{Kendine Referanslı Yapılar}

Bir yapı, kendi tipinden bir \emph{işaretçi} içerebilir. Bağlı liste, ağaç
gibi veri yapılarının temelidir.

\begin{lstlisting}
struct Dugum {
    int veri;
    struct Dugum *sonraki;   // kendine referans (isaretci olmali!)
};
\end{lstlisting}

\uyari{Yapı kendi tipini \emph{değer olarak} içeremez (\code{struct Dugum
sonraki;}) --- bu sonsuz boyut demektir. Yalnızca işaretçi (\code{*}) geçerlidir.}

\gsub{Esnek Dizi Üyesi (Flexible Array Member, C99)}

Yapının sonundaki boyutsuz dizi, yapı + veriyi \emph{tek} bellek bloğunda
tutmayı sağlar.

\begin{lstlisting}
#include <stdlib.h>
#include <string.h>
struct Metin {
    size_t uzunluk;
    char   veri[];        // esnek dizi uyesi (boyutsuz, en sonda)
};
struct Metin *metin_olustur(const char *s) {
    size_t n = strlen(s);
    struct Metin *m = malloc(sizeof(struct Metin) + n + 1); // yapi + veri
    if (!m) return NULL;
    m->uzunluk = n;
    strcpy(m->veri, s);
    return m;             // tek free ile hepsi temizlenir
}
\end{lstlisting}

% =====================================================================
\gsec{Birlikler (union) ve Enum}
% =====================================================================

\gsub{union: Aynı Belleği Paylaşan Üyeler}

Bir \code{union}, tüm üyelerini \emph{aynı bellek adresinde} tutar. Boyutu,
en büyük üyesi kadardır. Aynı anda yalnızca \emph{bir} üye anlamlıdır.

\begin{lstlisting}
#include <stdio.h>
union Deger {
    int    i;
    float  f;
    char   baytlar[4];
};
int main(void) {
    union Deger d;
    d.i = 1;                 // simdi 'i' anlamli
    printf("i = %d\n", d.i); // 1
    d.f = 3.14f;             // ayni bellege yazdi; 'i' artik gecersiz
    printf("f = %g\n", d.f); // 3.14
    printf("boyut = %zu\n", sizeof(union Deger)); // 4 (en buyuk uye)
    return 0;
}
\end{lstlisting}

\gsub{Type Punning: Bir Float'ın Bitlerini Görme}

Union, bir değeri farklı tip gözüyle okumaya (type punning) izin verir ---
IEEE 754 float'ın ham bitlerini incelemek için kullanışlıdır.

\begin{lstlisting}
#include <stdio.h>
#include <stdint.h>
int main(void) {
    union { float f; uint32_t u; } x;
    x.f = 1.0f;
    printf("1.0f'in bitleri: 0x%08X\n", x.u); // 0x3F800000
    return 0;
}
\end{lstlisting}

\bilgi{Endianness'i (bayt sırası) test etmenin klasik yolu da union'dır:
küçük-endian makinede \code{union\{int i; char c;\} x; x.i=1;} sonucunda
\code{x.c == 1} olur.}

\gsub{Tagged Union (Etiketli Birleşim)}

Union'ın hangi üyesinin geçerli olduğunu bilmek için genelde bir \code{enum}
"etiket" ile birlikte, bir \code{struct} içinde kullanılır. Bu desen, C'de
tip-güvenli değişken (variant) oluşturmanın standart yoludur.

\begin{lstlisting}
#include <stdio.h>
typedef enum { TIP_INT, TIP_FLOAT, TIP_METIN } Tip;

typedef struct {
    Tip tip;                    // hangi uye gecerli?
    union {
        int         i;
        float       f;
        const char *s;
    } u;
} Deger;

void yazdir(Deger d) {
    switch (d.tip) {
        case TIP_INT:   printf("int: %d\n",   d.u.i); break;
        case TIP_FLOAT: printf("float: %g\n", d.u.f); break;
        case TIP_METIN: printf("metin: %s\n", d.u.s); break;
    }
}
int main(void) {
    Deger a = { .tip = TIP_INT,   .u.i = 42 };
    Deger b = { .tip = TIP_METIN, .u.s = "merhaba" };
    yazdir(a); yazdir(b);
    return 0;
}
\end{lstlisting}

\gsub{enum: İsimli Sabitler}

\begin{lstlisting}
#include <stdio.h>
enum Renk { KIRMIZI, YESIL = 5, MAVI };  // 0, 5, 6

enum Gun { PZT = 1, SAL, CAR, PER, CUM }; // 1,2,3,4,5

int main(void) {
    enum Renk r = MAVI;
    printf("%d\n", r);        // 6
    printf("%d\n", CUM);      // 5
    return 0;
}
\end{lstlisting}

\gsub{Bit Alanları (Bit Fields)}

Yapı üyelerine bit düzeyinde boyut vererek belleği son derece verimli
kullanmayı sağlar --- donanım register'ları ve protokol başlıkları için ideal.

\begin{lstlisting}
#include <stdio.h>
struct Bayraklar {
    unsigned aktif   : 1;   // 1 bit
    unsigned gizli   : 1;   // 1 bit
    unsigned oncelik : 3;   // 3 bit (0-7)
    unsigned         : 3;   // 3 bit isimsiz dolgu
};                          // toplam 1 bayta sigar
int main(void) {
    struct Bayraklar b = {0};
    b.aktif = 1;
    b.oncelik = 5;
    printf("aktif=%u oncelik=%u boyut=%zu\n",
           b.aktif, b.oncelik, sizeof(b));  // aktif=1 oncelik=5 boyut=4
    return 0;
}
\end{lstlisting}

\gsub{Bit İşlemleri ve Bayrak Maskeleri}

Bit alanlarına alternatif olarak, bayrakları tek bir tamsayıda maskelerle de
saklayabilirsin --- taşınabilir ve hızlıdır.

\begin{lstlisting}
#include <stdio.h>
#define BAYRAK_OKU    (1u << 0)   // 0001
#define BAYRAK_YAZ    (1u << 1)   // 0010
#define BAYRAK_CALIS  (1u << 2)   // 0100

int main(void) {
    unsigned izin = 0;
    izin |=  BAYRAK_OKU | BAYRAK_YAZ;  // bit ekle (set)
    izin &= ~BAYRAK_YAZ;               // bit temizle (clear)
    izin ^=  BAYRAK_CALIS;             // bit degistir (toggle)

    if (izin & BAYRAK_OKU)             // bit test et
        printf("okuma izni var\n");
    printf("izin = %u\n", izin);       // 5
    return 0;
}
\end{lstlisting}

\begin{center}
\begin{tabularx}{\textwidth}{@{}>{\ttfamily}c l >{\raggedright\arraybackslash}X@{}}
\toprule
\rowcolor{acikMavi}\color{anaMavi}\textbf{Op} & \color{anaMavi}\textbf{Ad} & \color{anaMavi}\textbf{Kullanım}\\
\midrule
\& & AND & Bit test etme, maskeleme.\\
\rowcolor{acikGri} | & OR & Bit ekleme (set).\\
\^{} & XOR & Bit değiştirme (toggle), takas.\\
\rowcolor{acikGri} \textasciitilde & NOT & Tüm bitleri terse çevirme.\\
<< >> & Kaydırma & 2'nin kuvvetiyle hızlı çarpma/bölme.\\
\bottomrule
\end{tabularx}
\end{center}

% =====================================================================
\gsec{Fonksiyon İşaretçileri ve Callback'ler}
% =====================================================================

Fonksiyonların da bir adresi vardır; bu adresi bir işaretçide tutup daha
sonra çağırabiliriz. Bu, C'de \emph{davranışı veri gibi geçirmenin} yoludur ---
callback'ler, dispatch tabloları ve jenerik algoritmaların temeli.

\gsub{Söz Dizimi}

\begin{lstlisting}
// int alan, int donduren bir fonksiyonu gosteren isaretci:
int (*fp)(int);

// DIKKAT parantezler: int *fp(int) -> "int* donduren fonksiyon" (farkli!)
\end{lstlisting}

\begin{lstlisting}
#include <stdio.h>
int kare(int x)   { return x * x; }
int kup (int x)   { return x * x * x; }

int main(void) {
    int (*islem)(int);        // fonksiyon isaretcisi

    islem = kare;             // & opsiyonel: islem = &kare;
    printf("%d\n", islem(5)); // 25  (*islem)(5) ile ayni

    islem = kup;
    printf("%d\n", islem(3)); // 27
    return 0;
}
\end{lstlisting}

\gsub{Callback: qsort ile Sıralama}

Standart kütüphane \code{qsort}, karşılaştırmayı bir callback ile alır ---
böylece herhangi bir tipi, herhangi bir ölçüte göre sıralar.

\begin{lstlisting}
#include <stdio.h>
#include <stdlib.h>

// qsort'un istedigi imza: iki void* alir, int dondurur
int kucukten(const void *a, const void *b) {
    int x = *(const int *)a, y = *(const int *)b;
    return (x > y) - (x < y);   // tasmadan guvenli karsilastirma
}
int buyukten(const void *a, const void *b) {
    return kucukten(b, a);      // sirayi ters cevir
}

int main(void) {
    int a[] = {5, 2, 8, 1, 9, 3};
    size_t n = sizeof(a) / sizeof(a[0]);

    qsort(a, n, sizeof(int), kucukten);
    for (size_t i = 0; i < n; i++) printf("%d ", a[i]); // 1 2 3 5 8 9
    printf("\n");

    qsort(a, n, sizeof(int), buyukten);
    for (size_t i = 0; i < n; i++) printf("%d ", a[i]); // 9 8 5 3 2 1
    printf("\n");
    return 0;
}
\end{lstlisting}

\gsub{Fonksiyon İşaretçisi Dizisi (Dispatch Table)}

Uzun \code{switch}/\code{if} zincirleri yerine, fonksiyon işaretçilerinden
oluşan bir tablo ile O(1) yönlendirme yapılır.

\begin{lstlisting}
#include <stdio.h>
double topla(double a, double b)  { return a + b; }
double cikar(double a, double b)  { return a - b; }
double carp (double a, double b)  { return a * b; }
double bol  (double a, double b)  { return b ? a / b : 0; }

int main(void) {
    // Islem karakterini fonksiyona esleyen tablo
    struct { char op; double (*fn)(double, double); } tablo[] = {
        {'+', topla}, {'-', cikar}, {'*', carp}, {'/', bol},
    };
    char op = '*'; double a = 6, b = 7;
    for (size_t i = 0; i < sizeof(tablo)/sizeof(tablo[0]); i++)
        if (tablo[i].op == op) {
            printf("%g\n", tablo[i].fn(a, b));  // 42
            break;
        }
    return 0;
}
\end{lstlisting}

\gsub{typedef ile Okunabilir Fonksiyon İşaretçisi}

Fonksiyon işaretçisi tipleri hızla okunmaz hale gelir; \code{typedef} bunu
temizler.

\begin{lstlisting}
#include <stdio.h>
typedef int (*Islem)(int, int);   // artik 'Islem' bir tip

int topla(int a, int b) { return a + b; }

void uygula(Islem f, int x, int y) {   // parametre olarak temiz
    printf("%d\n", f(x, y));
}
int main(void) {
    Islem op = topla;               // okunakli bildirim
    uygula(op, 3, 4);               // 7
    return 0;
}
\end{lstlisting}

\ipucu{Fonksiyon işaretçileri jenerik programlamanın kalbidir: \code{qsort},
\code{bsearch}, \code{pthread\_create}, sinyal handler'ları ve olay
sistemleri hep callback alır.}

% =====================================================================
\gsec{Değişken Argümanlı Fonksiyonlar (Variadic)}
% =====================================================================

\code{printf} gibi değişken sayıda argüman alan fonksiyonlar
\code{<stdarg.h>} ile yazılır.

\begin{lstlisting}
#include <stdio.h>
#include <stdarg.h>

// Ilk parametre sayiyi verir; sonra o kadar int okunur.
int topla(int adet, ...) {
    va_list ap;
    va_start(ap, adet);          // 'adet'ten SONRAKI argumanlar
    int toplam = 0;
    for (int i = 0; i < adet; i++)
        toplam += va_arg(ap, int); // sirayla oku (tip belirt)
    va_end(ap);                    // temizle
    return toplam;
}
int main(void) {
    printf("%d\n", topla(3, 10, 20, 30));      // 60
    printf("%d\n", topla(5, 1, 2, 3, 4, 5));   // 15
    return 0;
}
\end{lstlisting}

\uyari{Variadic fonksiyonlar tip güvenli \emph{değildir}: \code{va\_arg}'a
yanlış tip verirsen tanımsız davranış olur. Argüman sayısını/tiplerini bir
biçim dizesi (printf gibi) ya da bir sayaç ile belirtmelisin. Sonlandırıcı
olarak \code{NULL} da kullanılabilir.}

% =====================================================================
\gsec{Önişlemci ve Makrolar}
% =====================================================================

Önişlemci, derlemeden önce metinsel dönüşümler yapar. Güçlüdür ama tuzaklıdır.

\gsub{Nesne ve Fonksiyon Benzeri Makrolar}

\begin{lstlisting}
#define PI 3.14159                     // nesne benzeri
#define KARE(x) ((x) * (x))            // fonksiyon benzeri
#define MAX(a, b) ((a) > (b) ? (a) : (b))
\end{lstlisting}

\uyari{Makro argümanlarını ve tüm ifadeyi \emph{parantezle}. \code{\#define
KARE(x) x*x} yazarsan \code{KARE(1+2)} $\to$ \code{1+2*1+2 = 5} olur (9 değil).
Ayrıca \code{MAX(i++, j++)} argümanı iki kez değerlendirir --- yan etkili
argüman verme.}

\gsub{Koşullu Derleme ve Include Guard}

\begin{lstlisting}
#ifndef BASLIK_H          // include guard: cift dahil etmeyi onler
#define BASLIK_H
/* ... bildirimler ... */
#endif

#ifdef DEBUG              // yalnizca -DDEBUG ile derlenince
    #define LOG(msg) fprintf(stderr, "[LOG] %s\n", msg)
#else
    #define LOG(msg) ((void)0)   // hicbir sey yapmaz
#endif

#if defined(__linux__)
    /* Linux'a ozgu kod */
#elif defined(_WIN32)
    /* Windows'a ozgu kod */
#endif
\end{lstlisting}

\gsub{Stringize (\#) ve Token Paste (\#\#)}

\begin{lstlisting}
#include <stdio.h>
#define STR(x)   #x            // argumani stringe cevirir
#define BIRLES(a, b) a##b      // iki token'i birlestirir

int main(void) {
    printf("%s\n", STR(merhaba dunya)); // "merhaba dunya"
    int xy = 5;
    printf("%d\n", BIRLES(x, y));       // xy -> 5
    return 0;
}
\end{lstlisting}

\gsub{Faydalı Hazır Makrolar}

\begin{center}
\begin{tabularx}{\textwidth}{@{}>{\ttfamily\color{anaMavi}}l >{\raggedright\arraybackslash}X@{}}
\toprule
\rowcolor{acikMavi}\color{anaMavi}\textbf{Makro} & \color{anaMavi}\textbf{Değer}\\
\midrule
\_\_FILE\_\_ & Kaynak dosya adı (string).\\
\rowcolor{acikGri} \_\_LINE\_\_ & Bulunulan satır numarası.\\
\_\_func\_\_ & İçinde bulunulan fonksiyonun adı (C99).\\
\rowcolor{acikGri} \_\_DATE\_\_ / \_\_TIME\_\_ & Derleme tarihi/saati.\\
\bottomrule
\end{tabularx}
\end{center}

\ornek[Hata ayıklama makrosu]{Dosya, satır ve fonksiyon bilgisiyle log.}
\begin{lstlisting}
#include <stdio.h>
#define HATA(fmt, ...) \
    fprintf(stderr, "%s:%d (%s): " fmt "\n", \
            __FILE__, __LINE__, __func__, ##__VA_ARGS__)

int main(void) {
    HATA("kod = %d", 42);   // main.c:8 (main): kod = 42
    return 0;
}
\end{lstlisting}

\gsub{X-Macros: Tek Kaynaktan Çoklu Üretim}

Aynı listeyi hem enum hem string dizisi olarak üretmek gibi tekrarları,
X-macro deseni tek bir yerden yönetir.

\begin{lstlisting}
#include <stdio.h>
#define RENKLER  \
    X(KIRMIZI)   \
    X(YESIL)     \
    X(MAVI)

typedef enum {
    #define X(ad) ad,
    RENKLER
    #undef X
} Renk;

static const char *renk_adi[] = {
    #define X(ad) #ad,
    RENKLER
    #undef X
};
int main(void) {
    Renk r = YESIL;
    printf("%d -> %s\n", r, renk_adi[r]);  // 1 -> YESIL
    return 0;
}
\end{lstlisting}

% =====================================================================
\gsec{Jenerik Programlama Teknikleri}
% =====================================================================

C'nin şablonları yoktur, ama \code{void *}, makrolar ve C11'in \code{\_Generic}
özelliğiyle jenerik kod yazılabilir.

\gsub{void * ile Jenerik Kap}

\begin{lstlisting}
#include <stdio.h>
#include <string.h>
// Herhangi bir tipteki elemani genel takas (swap):
void takas(void *a, void *b, size_t boyut) {
    unsigned char tmp[boyut];          // VLA (C99)
    memcpy(tmp, a,   boyut);
    memcpy(a,   b,   boyut);
    memcpy(b,   tmp, boyut);
}
int main(void) {
    int x = 1, y = 2;
    takas(&x, &y, sizeof(int));
    printf("%d %d\n", x, y);           // 2 1

    double p = 1.5, q = 3.5;
    takas(&p, &q, sizeof(double));
    printf("%g %g\n", p, q);           // 3.5 1.5
    return 0;
}
\end{lstlisting}

\gsub{\_Generic ile Tipe Göre Seçim (C11)}

\begin{lstlisting}
#include <stdio.h>
// Argumanin tipine gore farkli isim uretir
#define TIP_ADI(x) _Generic((x), \
    int:    "int",   \
    double: "double",\
    char:   "char",  \
    char *: "string",\
    default: "bilinmeyen")

int main(void) {
    printf("%s\n", TIP_ADI(42));      // int
    printf("%s\n", TIP_ADI(3.14));    // double
    printf("%s\n", TIP_ADI("selam")); // string
    return 0;
}
\end{lstlisting}

\bilgi{\code{\_Generic} derleme anında çalışır ve tip-güvenli "aşırı yükleme"
(overloading) benzeri davranış sağlar. Standart kütüphanedeki \code{<tgmath.h>}
tam olarak bunu kullanır.}

% =====================================================================
\gsec{Uygulama: İşaretçilerle Veri Yapıları}
% =====================================================================

Öğrenilenleri birleştiren en iyi alıştırma, veri yapılarını sıfırdan yazmaktır.

\gsub{Tekli Bağlı Liste}

\begin{lstlisting}
#include <stdio.h>
#include <stdlib.h>

typedef struct Dugum {
    int veri;
    struct Dugum *sonraki;
} Dugum;

// Basa ekle: O(1). Liste bası cift-isaretci ile guncellenir.
void basa_ekle(Dugum **bas, int veri) {
    Dugum *yeni = malloc(sizeof(Dugum));
    if (!yeni) return;
    yeni->veri = veri;
    yeni->sonraki = *bas;    // yeni dugum eski basi gosterir
    *bas = yeni;             // bas artik yeni dugum
}

void yazdir(const Dugum *bas) {
    for (const Dugum *p = bas; p; p = p->sonraki)
        printf("%d -> ", p->veri);
    printf("NULL\n");
}

// Tum listeyi serbest birak (bellek sizintisini onle)
void yok_et(Dugum *bas) {
    while (bas) {
        Dugum *sonraki = bas->sonraki; // ONCE sonrakini sakla
        free(bas);                     // sonra bu dugumu birak
        bas = sonraki;
    }
}

int main(void) {
    Dugum *liste = NULL;
    basa_ekle(&liste, 3);
    basa_ekle(&liste, 2);
    basa_ekle(&liste, 1);
    yazdir(liste);           // 1 -> 2 -> 3 -> NULL
    yok_et(liste);
    return 0;
}
\end{lstlisting}

\uyari{Düğümü \code{free} etmeden \emph{önce} \code{sonraki}'yi sakla. Ters
sıra yaparsan (önce free) asılı işaretçiye erişip çökeresin.}

\gsub{Dinamik Yığın (Stack) --- realloc ile Büyüyen Dizi}

\begin{lstlisting}
#include <stdio.h>
#include <stdlib.h>

typedef struct {
    int   *veri;
    size_t adet;      // dolu eleman
    size_t kapasite;  // ayrilmis eleman
} Yigin;

int yigin_it(Yigin *y, int deger) {
    if (y->adet == y->kapasite) {                 // dolduysa buyut
        size_t yeni_kap = y->kapasite ? y->kapasite * 2 : 4;
        int *yeni = realloc(y->veri, yeni_kap * sizeof(int));
        if (!yeni) return -1;                     // hata: eski gecerli
        y->veri = yeni;
        y->kapasite = yeni_kap;
    }
    y->veri[y->adet++] = deger;
    return 0;
}
int yigin_cek(Yigin *y, int *out) {
    if (y->adet == 0) return -1;                  // bos
    *out = y->veri[--y->adet];
    return 0;
}
int main(void) {
    Yigin y = {0};
    for (int i = 1; i <= 6; i++) yigin_it(&y, i * 10);
    int v;
    while (yigin_cek(&y, &v) == 0) printf("%d ", v); // 60 50 40 30 20 10
    printf("\n");
    free(y.veri);
    return 0;
}
\end{lstlisting}

\ipucu{Bu "büyüyen dizi" (dynamic array / vector) deseni --- kapasiteyi
ikiye katlama --- amortize edilmiş O(1) ekleme sağlar ve gerçek dünyada
en sık kullanılan veri yapısıdır.}

% #####################################################################
\gpart{2}{Linux Sistem Programlama}
% #####################################################################

% =====================================================================
\gsec{Sistem Çağrıları, Kernel ve errno}
% =====================================================================

Sistem programlama, programın işletim sistemi çekirdeğiyle (kernel) doğrudan
konuşmasıdır: dosya açmak, süreç oluşturmak, ağ üzerinden veri göndermek.
Bu işlemler \emph{sistem çağrıları} (system calls) ile yapılır.

\gsub{Kullanıcı Alanı ve Çekirdek Alanı}

Program normalde \emph{kullanıcı alanında} (user space) çalışır; donanıma
doğrudan erişemez. Bir kaynağa (dosya, ağ, bellek) ihtiyaç duyduğunda bir
sistem çağrısı yaparak \emph{çekirdek alanına} (kernel space) geçer. Çekirdek
işi güvenli biçimde yapıp sonucu döndürür.

\begin{center}
\begin{tabularx}{\textwidth}{@{}>{\raggedright\arraybackslash}X >{\raggedright\arraybackslash}X@{}}
\toprule
\rowcolor{acikMavi}\color{anaMavi}\textbf{Kütüphane Fonksiyonu (libc)} & \color{anaMavi}\textbf{Sistem Çağrısı (kernel)}\\
\midrule
\code{fopen, fread, printf} (tamponlu) & \code{open, read, write} (tamponsuz)\\
\rowcolor{acikGri}\code{malloc} & \code{brk, sbrk, mmap}\\
\code{system} & \code{fork + execve + wait}\\
\bottomrule
\end{tabularx}
\end{center}

\bilgi{\code{man 2 <ad>} bir \emph{sistem çağrısını} anlatır (bölüm 2),
\code{man 3 <ad>} ise bir \emph{kütüphane fonksiyonunu} (bölüm 3).
Örnek: \code{man 2 open}, \code{man 3 printf}.}

\gsub{Hata Yönetimi: errno}

Çoğu sistem çağrısı başarısızlıkta \code{-1} (ya da \code{NULL}) döndürür ve
küresel \code{errno} değişkenine hata kodunu yazar. \code{perror} ve
\code{strerror} bunu okunur mesaja çevirir.

\begin{lstlisting}
#include <stdio.h>
#include <string.h>
#include <errno.h>
#include <fcntl.h>
#include <unistd.h>

int main(void) {
    int fd = open("/olmayan/dosya", O_RDONLY);
    if (fd == -1) {                       // hata olustu
        // errno'yu hemen kullan (sonraki cagri ustune yazabilir)
        fprintf(stderr, "Hata kodu: %d\n", errno);
        perror("open");                   // "open: No such file or directory"
        fprintf(stderr, "strerror: %s\n", strerror(errno));
        return 1;
    }
    close(fd);
    return 0;
}
\end{lstlisting}

\uyari{\code{errno} yalnızca bir çağrı \emph{başarısız olduğunda} anlamlıdır;
başarıda sıfırlanmaz. Bu yüzden önce dönüş değerini kontrol et, \emph{sonra}
errno'ya bak. Ayrıca errno'yu araya başka bir çağrı girmeden hemen kullan.}

% =====================================================================
\gsec{Düşük Seviye Dosya G/Ç}
% =====================================================================

\code{<stdio.h>}'nin tamponlu \code{FILE*} akışlarının altında, çekirdeğin
\emph{dosya betimleyicileri} (file descriptor --- fd) yatar. Bunlar küçük
tamsayılardır: 0 = stdin, 1 = stdout, 2 = stderr.

\gsub{open, read, write, close}

\begin{center}
\begin{tabularx}{\textwidth}{@{}>{\ttfamily\color{anaMavi}}l >{\raggedright\arraybackslash}X@{}}
\toprule
\rowcolor{acikMavi}\color{anaMavi}\textbf{Çağrı} & \color{anaMavi}\textbf{Görev}\\
\midrule
open(yol, bayrak, mod) & Dosya açar/oluşturur; fd döndürür.\\
\rowcolor{acikGri} read(fd, tampon, n) & En çok n bayt okur; okunan sayıyı döndürür.\\
write(fd, tampon, n) & En çok n bayt yazar; yazılan sayıyı döndürür.\\
\rowcolor{acikGri} lseek(fd, ofset, nereden) & Dosya konumunu değiştirir.\\
close(fd) & Betimleyiciyi kapatır.\\
\bottomrule
\end{tabularx}
\end{center}

\gsubsub{open Bayrakları}

\begin{center}
\begin{tabularx}{\textwidth}{@{}>{\ttfamily\color{anaMavi}}l >{\raggedright\arraybackslash}X@{}}
\toprule
\rowcolor{acikMavi}\color{anaMavi}\textbf{Bayrak} & \color{anaMavi}\textbf{Anlam}\\
\midrule
O\_RDONLY / O\_WRONLY / O\_RDWR & Salt oku / salt yaz / oku-yaz.\\
\rowcolor{acikGri} O\_CREAT & Yoksa oluştur (3. argüman: izin modu, örn. 0644).\\
O\_TRUNC & Varsa içeriği sil (sıfırla).\\
\rowcolor{acikGri} O\_APPEND & Sona ekle (her yazma dosya sonuna).\\
O\_EXCL & O\_CREAT ile: dosya zaten varsa hata ver.\\
\bottomrule
\end{tabularx}
\end{center}

\ornek[Dosya kopyalama (cp'nin özü)]{Bir dosyayı bayt bayt okuyup diğerine
yazan tam program. Kısmi okuma/yazmaya ve sinyal kesintilerine dikkat edilir.}
\begin{lstlisting}
#include <fcntl.h>
#include <unistd.h>
#include <stdio.h>
#include <errno.h>

int main(int argc, char *argv[]) {
    if (argc != 3) {
        fprintf(stderr, "Kullanim: %s kaynak hedef\n", argv[0]);
        return 1;
    }
    int kaynak = open(argv[1], O_RDONLY);
    if (kaynak == -1) { perror("open kaynak"); return 1; }

    // 0644 = kullanici oku/yaz, grup+digerleri oku
    int hedef = open(argv[2], O_WRONLY | O_CREAT | O_TRUNC, 0644);
    if (hedef == -1) { perror("open hedef"); close(kaynak); return 1; }

    char tampon[4096];
    ssize_t okunan;
    while ((okunan = read(kaynak, tampon, sizeof(tampon))) > 0) {
        ssize_t yazilan = 0;
        while (yazilan < okunan) {                 // kismi yazmayi ele al
            ssize_t r = write(hedef, tampon + yazilan, okunan - yazilan);
            if (r == -1) {
                if (errno == EINTR) continue;      // sinyal kesti; tekrar dene
                perror("write"); close(kaynak); close(hedef); return 1;
            }
            yazilan += r;
        }
    }
    if (okunan == -1) perror("read");

    close(kaynak);
    close(hedef);
    return 0;
}
\end{lstlisting}

\uyari{\code{read}/\code{write} istediğin kadar bayt işlemeyebilir (kısmi
transfer). \code{write} özellikle boru ve sokette daha az yazabilir. Bu yüzden
\emph{her zaman} döngüyle, dönüş değerini kontrol ederek yaz. \code{EINTR}
(sinyalle kesilme) durumunda çağrıyı tekrarla.}

\gsub{lseek ve Dosya Konumu}

\begin{lstlisting}
#include <fcntl.h>
#include <unistd.h>
#include <stdio.h>
int main(void) {
    int fd = open("veri.bin", O_RDWR | O_CREAT, 0644);

    write(fd, "MERHABA", 7);
    lseek(fd, 0, SEEK_SET);          // basa don
    char buf[8] = {0};
    read(fd, buf, 7);
    printf("%s\n", buf);             // MERHABA

    // Seyrek dosya (sparse file): ileri atla, delik olusur
    lseek(fd, 1000, SEEK_END);       // sondan 1000 bayt ileri
    write(fd, "SON", 3);

    off_t boyut = lseek(fd, 0, SEEK_END); // dosya boyutu
    printf("boyut = %ld\n", (long)boyut);

    close(fd);
    return 0;
}
\end{lstlisting}

\notk[SEEK sabitleri]{\code{SEEK\_SET} = dosya başından, \code{SEEK\_CUR} =
mevcut konumdan, \code{SEEK\_END} = dosya sonundan itibaren ofset.}

\gsub{dup, dup2 ve Betimleyici Yönlendirme}

\code{dup2} bir betimleyiciyi başka bir numaraya kopyalar --- kabuğun
\code{>} yönlendirmesinin ve boru bağlamanın temelidir.

\begin{lstlisting}
#include <fcntl.h>
#include <unistd.h>
int main(void) {
    int fd = open("cikti.txt", O_WRONLY | O_CREAT | O_TRUNC, 0644);
    dup2(fd, STDOUT_FILENO);   // stdout artik dosyaya gider
    close(fd);

    // Bu andan sonra tum stdout ciktisi cikti.txt'ye yazilir
    write(STDOUT_FILENO, "Bu satir dosyaya gitti\n", 23);
    return 0;
}
\end{lstlisting}

\gsub{Betimleyici mi, Akış mı?}

\begin{center}
\begin{tabularx}{\textwidth}{@{}>{\raggedright\arraybackslash}X >{\raggedright\arraybackslash}X@{}}
\toprule
\rowcolor{acikMavi}\color{anaMavi}\textbf{fd (open/read/write)} & \color{anaMavi}\textbf{FILE* (fopen/fread)}\\
\midrule
Ham sistem çağrısı, tamponsuz & Kullanıcı alanında tamponlu (hızlı)\\
\rowcolor{acikGri} Küçük tamsayı & \code{FILE} yapısına işaretçi\\
Sokete/boruya doğrudan uyar & Satır/tam tamponlama; \code{fflush} gerekir\\
\bottomrule
\end{tabularx}
\end{center}

% =====================================================================
\gsec{İleri Dosya G/Ç}
% =====================================================================

Temel \code{open/read/write} üzerine, gerçek uygulamaların ihtiyaç duyduğu
konumsal G/Ç, dağınık G/Ç, dosya kilitleme, kalıcılık (durability), geçici
dosyalar ve verimli kopyalama teknikleri.

\gsub{Konumsal G/Ç: pread ve pwrite}

\code{pread}/\code{pwrite}, dosya konumunu (offset) \emph{değiştirmeden}
belirli bir konumdan okur/yazar. \code{lseek + read}'in aksine atomiktir ---
iş parçacıklarının aynı dosyayı yarışmadan kullanmasını sağlar.

\begin{lstlisting}
#include <fcntl.h>
#include <unistd.h>
#include <stdio.h>
int main(void) {
    int fd = open("veri.bin", O_RDWR | O_CREAT, 0644);

    char kayit[16] = "KAYIT-0000000000";
    pwrite(fd, kayit, 16, 0);        // 0. konuma yaz (offset degismez)
    pwrite(fd, kayit, 16, 1600);     // 100. kaydin yerine yaz

    char tampon[16];
    pread(fd, tampon, 16, 1600);     // dosya konumundan BAGIMSIZ oku
    printf("%.16s\n", tampon);

    close(fd);
    return 0;
}
\end{lstlisting}

\bilgi{Çok iş parçacıklı bir programda tek bir açık dosyaya paralel erişim
için \code{pread}/\code{pwrite} şarttır: dosya konumu paylaşıldığından
\code{lseek}+\code{read} ikilisi yarış durumu yaratır.}

\gsub{Dağınık/Toplu G/Ç: readv ve writev (Scatter-Gather)}

Tek sistem çağrısıyla birden çok tampona okur / birden çok tampondan yazar.
Başlık + gövdeyi ayrı tamponlardan tek atomik yazmayla göndermek için idealdir.

\begin{lstlisting}
#include <sys/uio.h>
#include <unistd.h>
#include <string.h>
int main(void) {
    const char *baslik = "[BASLIK] ";
    const char *govde   = "mesaj govdesi\n";

    struct iovec parcalar[2];
    parcalar[0].iov_base = (void *)baslik;
    parcalar[0].iov_len  = strlen(baslik);
    parcalar[1].iov_base = (void *)govde;
    parcalar[1].iov_len  = strlen(govde);

    // Iki tampon TEK writev cagrisiyla, sirayla, atomik yazilir
    writev(STDOUT_FILENO, parcalar, 2);   // [BASLIK] mesaj govdesi
    return 0;
}
\end{lstlisting}

\gsub{Dosya Bayraklarını Yönetme: fcntl}

\code{fcntl}, açık bir betimleyicinin özelliklerini çalışma anında değiştirir:
non-blocking moda geçme, \code{FD\_CLOEXEC} (exec'te otomatik kapanma) gibi.

\begin{lstlisting}
#include <fcntl.h>
#include <unistd.h>
int main(void) {
    int fd = open("dosya.txt", O_RDONLY);

    // Mevcut bayraklari al, O_NONBLOCK ekle, geri yaz
    int bayraklar = fcntl(fd, F_GETFL);
    fcntl(fd, F_SETFL, bayraklar | O_NONBLOCK);  // artik bloklamaz

    // exec sonrasi otomatik kapansin (betimleyici sizintisini onler)
    int fbayrak = fcntl(fd, F_GETFD);
    fcntl(fd, F_SETFD, fbayrak | FD_CLOEXEC);

    close(fd);
    return 0;
}
\end{lstlisting}

\uyari{Betimleyiciler \code{fork} sonrası çocukta açık kalır ve \code{exec}
sonrası da (varsayılan) açık kalır --- bu bir güvenlik/kaynak sızıntısıdır.
Hassas dosyaları \code{FD\_CLOEXEC} ile işaretle ya da \code{open}'da
\code{O\_CLOEXEC} bayrağını kullan.}

\gsub{Dosya Kilitleme (fcntl F\_SETLK)}

Birden çok süreç aynı dosyaya yazarken tutarlılık için öneri kilitleri
(advisory locks) kullanılır.

\begin{lstlisting}
#include <fcntl.h>
#include <unistd.h>
#include <stdio.h>
int main(void) {
    int fd = open("kayit.db", O_RDWR | O_CREAT, 0644);

    struct flock kilit = {0};
    kilit.l_type   = F_WRLCK;    // yazma kilidi (F_RDLCK: okuma, F_UNLCK: ac)
    kilit.l_whence = SEEK_SET;
    kilit.l_start  = 0;          // ofsetten...
    kilit.l_len    = 0;          // ...sona kadar (0 = tum dosya)

    if (fcntl(fd, F_SETLKW, &kilit) == -1) { // F_SETLKW: kilit bosalana kadar bekle
        perror("kilit"); return 1;
    }
    // --- kritik bolge: dosyaya guvenle yaz ---
    pwrite(fd, "guncel", 6, 0);

    kilit.l_type = F_UNLCK;
    fcntl(fd, F_SETLK, &kilit);  // kilidi birak
    close(fd);
    return 0;
}
\end{lstlisting}

\notk[F\_SETLK vs F\_SETLKW]{\code{F\_SETLK} kilit meşgulse hemen \code{-1}
döner (EAGAIN). \code{F\_SETLKW} ise kilit serbest kalana kadar \emph{bekler}
(W = wait). Basit dosya kilidi için \code{flock(fd, LOCK\_EX)} da kullanılabilir.}

\gsub{Kalıcılık: fsync, fdatasync}

\code{write} verinin diske yazıldığını \emph{garanti etmez} --- çekirdek onu
önbellekte tutar. Veritabanları ve dosya bütünlüğü için diske zorlamak gerekir.

\begin{lstlisting}
#include <fcntl.h>
#include <unistd.h>
int main(void) {
    int fd = open("onemli.log", O_WRONLY | O_CREAT | O_APPEND, 0644);
    write(fd, "kritik islem tamam\n", 19);
    fsync(fd);        // veri + metadata'yi diske zorla (guc kesintisine dayanikli)
    // fdatasync(fd); // yalnizca veri (boyut degismediyse daha hizli)
    close(fd);
    return 0;
}
\end{lstlisting}

\gsub{Güvenli Geçici Dosya: mkstemp}

\begin{lstlisting}
#include <stdlib.h>
#include <unistd.h>
#include <stdio.h>
int main(void) {
    char sablon[] = "/tmp/uygXXXXXX";   // son 6 karakter X olmali
    int fd = mkstemp(sablon);           // benzersiz ad uretir + acar (guvenli)
    if (fd == -1) { perror("mkstemp"); return 1; }
    printf("Gecici dosya: %s\n", sablon);

    write(fd, "gecici veri", 11);
    // Genellikle isim hemen silinir; fd acik kaldigi surece dosya yasar
    unlink(sablon);                     // program bitince tamamen yok olur
    close(fd);
    return 0;
}
\end{lstlisting}

\uyari{Asla \code{tmpnam}/\code{mktemp} kullanma --- ad üretimiyle açma
arasında yarış (TOCTOU) güvenlik açığı vardır. \code{mkstemp} adı üretir ve
atomik olarak açar.}

\gsub{Verimli Kopyalama: sendfile}

\code{sendfile}, veriyi kullanıcı alanına taşımadan çekirdek içinde bir fd'den
diğerine kopyalar --- dosya sunucularında (statik içerik) çok hızlıdır.

\begin{lstlisting}
#include <sys/sendfile.h>
#include <sys/stat.h>
#include <fcntl.h>
#include <unistd.h>
int main(int argc, char *argv[]) {
    int giris = open(argv[1], O_RDONLY);
    int cikis = open(argv[2], O_WRONLY | O_CREAT | O_TRUNC, 0644);
    struct stat st; fstat(giris, &st);

    off_t ofset = 0;
    // Tum dosyayi cekirdek icinde kopyala (read/write dongusu YOK)
    sendfile(cikis, giris, &ofset, st.st_size);

    close(giris); close(cikis);
    return 0;
}
\end{lstlisting}

\gsub{stdio Tamponlaması: setvbuf}

\code{FILE*} akışları kullanıcı alanında tamponlanır. Tamponlama modunu
kontrol etmek, performans ve etkileşimli çıktı için önemlidir.

\begin{center}
\begin{tabularx}{\textwidth}{@{}>{\ttfamily\color{anaMavi}}l >{\raggedright\arraybackslash}X@{}}
\toprule
\rowcolor{acikMavi}\color{anaMavi}\textbf{Mod} & \color{anaMavi}\textbf{Davranış}\\
\midrule
\_IOFBF & Tam tamponlama: tampon dolunca yazılır (dosyalarda varsayılan).\\
\rowcolor{acikGri} \_IOLBF & Satır tamponlama: \code{\textbackslash n} görünce yazılır (terminalde varsayılan).\\
\_IONBF & Tamponsuz: her yazma anında gider (stderr).\\
\bottomrule
\end{tabularx}
\end{center}

\begin{lstlisting}
#include <stdio.h>
int main(void) {
    // stdout'u tamponsuz yap (log ciktisi aninda gorunsun)
    setvbuf(stdout, NULL, _IONBF, 0);
    printf("Bu satir aninda gorunur");   // \n olmadan bile
    // Alternatif: fflush(stdout); ile tamponu elle bosalt
    return 0;
}
\end{lstlisting}

\ipucu{Program çökerse tamponlanmış \code{printf} çıktısı \emph{kaybolabilir}.
Hata ayıklarken ya \code{stderr} kullan (tamponsuz), ya her kritik noktada
\code{fflush(stdout)} çağır, ya da \code{setvbuf} ile tamponlamayı kapat.}

% =====================================================================
\gsec{Dosya Sistemi: stat, Dizinler, İzinler}
% =====================================================================

\gsub{stat ile Dosya Bilgisi}

\begin{lstlisting}
#include <sys/stat.h>
#include <stdio.h>
#include <time.h>
int main(int argc, char *argv[]) {
    if (argc != 2) { fprintf(stderr, "Kullanim: %s dosya\n", argv[0]); return 1; }
    struct stat st;
    if (stat(argv[1], &st) == -1) { perror("stat"); return 1; }

    printf("Boyut     : %ld bayt\n", (long)st.st_size);
    printf("Inode     : %ld\n", (long)st.st_ino);
    printf("Link sayisi: %ld\n", (long)st.st_nlink);
    printf("Izinler   : %o\n", st.st_mode & 0777);
    printf("Sahip UID : %d\n", st.st_uid);

    if (S_ISREG(st.st_mode))  printf("Tur: normal dosya\n");
    if (S_ISDIR(st.st_mode))  printf("Tur: dizin\n");
    if (S_ISLNK(st.st_mode))  printf("Tur: sembolik baglanti\n");

    printf("Son degisim: %s", ctime(&st.st_mtime));
    return 0;
}
\end{lstlisting}

\bilgi{\code{stat} sembolik bağlantıyı \emph{takip eder} (hedefi anlatır);
bağlantının kendisini incelemek için \code{lstat}, açık bir fd için ise
\code{fstat} kullanılır.}

\gsub{Dizin Okuma (opendir/readdir)}

\begin{lstlisting}
#include <dirent.h>
#include <stdio.h>
int main(int argc, char *argv[]) {
    const char *yol = (argc > 1) ? argv[1] : ".";
    DIR *d = opendir(yol);
    if (!d) { perror("opendir"); return 1; }

    struct dirent *girdi;
    while ((girdi = readdir(d)) != NULL)     // her cagri bir girdi
        printf("%s\n", girdi->d_name);        // . ve .. dahil

    closedir(d);
    return 0;
}
\end{lstlisting}

\gsub{İzinler: chmod, access, umask}

\begin{lstlisting}
#include <sys/stat.h>
#include <unistd.h>
int main(void) {
    chmod("betik.sh", 0755);        // rwxr-xr-x

    if (access("dosya.txt", R_OK) == 0)   // okuma izni var mi?
        /* okunabilir */;
    if (access("dosya.txt", F_OK) == 0)   // dosya var mi?
        /* mevcut */;
    return 0;
}
\end{lstlisting}

\notk[Sekizlik izin gösterimi]{İzinler üç haneli sekizlik sayıdır:
sahip/grup/diğerleri. Her hane oku(4)+yaz(2)+çalış(1) toplamıdır.
Örnek: 0644 = rw-r--r--, 0755 = rwxr-xr-x, 0600 = rw-------.}

\gsub{Dizin ve Bağlantı İşlemleri}

\begin{center}
\begin{tabularx}{\textwidth}{@{}>{\ttfamily\color{anaMavi}}l >{\raggedright\arraybackslash}X@{}}
\toprule
\rowcolor{acikMavi}\color{anaMavi}\textbf{Çağrı} & \color{anaMavi}\textbf{Görev}\\
\midrule
mkdir(yol, mod) & Dizin oluştur.\\
\rowcolor{acikGri} rmdir(yol) & Boş dizini sil.\\
unlink(yol) & Dosyayı (bağlantıyı) sil.\\
\rowcolor{acikGri} rename(eski, yeni) & Yeniden adlandır/taşı (atomik).\\
link(hedef, ad) & Sabit bağlantı (hard link) oluştur.\\
\rowcolor{acikGri} symlink(hedef, ad) & Sembolik bağlantı oluştur.\\
readlink(yol, buf, n) & Sembolik bağlantının hedefini oku.\\
\rowcolor{acikGri} chdir(yol) / getcwd(buf, n) & Çalışma dizinini değiştir / öğren.\\
\bottomrule
\end{tabularx}
\end{center}

\begin{lstlisting}
#include <sys/stat.h>
#include <unistd.h>
#include <stdio.h>
int main(void) {
    mkdir("test_dizin", 0755);              // dizin olustur

    int fd = open("test_dizin/a.txt", O_WRONLY | O_CREAT, 0644);
    write(fd, "veri", 4); close(fd);

    link("test_dizin/a.txt", "test_dizin/b.txt");   // sabit baglanti
    symlink("a.txt", "test_dizin/c.txt");           // sembolik baglanti
    rename("test_dizin/a.txt", "test_dizin/d.txt"); // yeniden adlandir

    char cwd[256];
    getcwd(cwd, sizeof cwd);                // mevcut dizin
    printf("Calisma dizini: %s\n", cwd);

    // Temizlik
    unlink("test_dizin/b.txt");
    unlink("test_dizin/c.txt");
    unlink("test_dizin/d.txt");
    rmdir("test_dizin");                    // dizin bosalinca silinebilir
    return 0;
}
\end{lstlisting}

\notk[Sabit vs Sembolik Bağlantı]{\textbf{Sabit bağlantı} aynı inode'a ikinci
bir isimdir; asıl silinse bile veri yaşar (link sayısı 0 olana dek).
\textbf{Sembolik bağlantı} başka bir \emph{yola} işaret eden küçük bir
dosyadır; hedef silinirse "kırık" (dangling) olur.}

\gsub{Dizin Ağacında Gezinme (nftw)}

\code{nftw}, bir dizin ağacını özyinelemeli olarak dolaşıp her girdi için bir
callback çağırır --- \code{find} komutunun temelidir.

\begin{lstlisting}
#define _XOPEN_SOURCE 500
#include <ftw.h>
#include <stdio.h>

// Her dosya/dizin icin cagrilir
int ziyaret(const char *yol, const struct stat *st, int tip, struct FTW *f) {
    (void)f;
    printf("%*s%s (%ld bayt)\n", f->level * 2, "", yol, (long)st->st_size);
    return 0;   // 0: devam et, sifir disi: durdur
}
int main(int argc, char *argv[]) {
    const char *kok = (argc > 1) ? argv[1] : ".";
    // 20: ayni anda acik fd siniri; FTW_PHYS: sembolik baglantilari izleme
    nftw(kok, ziyaret, 20, FTW_PHYS);
    return 0;
}
\end{lstlisting}

% =====================================================================
\gsec{G/Ç Çoğullama: select, poll, epoll}
% =====================================================================

Tek bir iş parçacığının \emph{birden çok} betimleyiciyi (soket, boru, terminal)
aynı anda izlemesini sağlar: "hangisi okumaya hazır?" Ağ sunucularının kalbidir.

\gsub{Neden Gerekli?}

Bloklayan \code{read} tek bir fd'de bekler; o gelene kadar diğerleri
işlenemez. G/Ç çoğullama, "bu fd'lerden herhangi biri hazır olunca beni
uyandır" der --- böylece tek iş parçacığı binlerce bağlantıyı yönetir.

\gsub{select}

\begin{lstlisting}
#include <sys/select.h>
#include <unistd.h>
#include <stdio.h>
int main(void) {
    fd_set okunacak;
    struct timeval zaman_asimi = { .tv_sec = 5, .tv_usec = 0 };

    FD_ZERO(&okunacak);              // kumeyi temizle
    FD_SET(STDIN_FILENO, &okunacak); // stdin'i izle

    // 1. arg: en buyuk fd + 1. Doner: hazir fd sayisi (0 = zaman asimi)
    int hazir = select(STDIN_FILENO + 1, &okunacak, NULL, NULL, &zaman_asimi);

    if (hazir == -1)      perror("select");
    else if (hazir == 0)  printf("5 saniye icinde giris yok.\n");
    else if (FD_ISSET(STDIN_FILENO, &okunacak)) {
        char buf[128];
        ssize_t n = read(STDIN_FILENO, buf, sizeof buf - 1);
        if (n > 0) { buf[n] = '\0'; printf("Girdi: %s", buf); }
    }
    return 0;
}
\end{lstlisting}

\uyari{\code{select}'in sınırları vardır: fd sayısı \code{FD\_SETSIZE} (genelde
1024) ile sınırlıdır ve her çağrıda küme yeniden kurulmalıdır. Çok sayıda
bağlantı için \code{poll} ya da \code{epoll} tercih edilir.}

\gsub{poll}

\begin{lstlisting}
#include <poll.h>
#include <unistd.h>
#include <stdio.h>
int main(void) {
    struct pollfd fds[1];
    fds[0].fd     = STDIN_FILENO;
    fds[0].events = POLLIN;          // okumaya hazir olmasini izle

    int hazir = poll(fds, 1, 5000);  // 5000 ms zaman asimi
    if (hazir > 0 && (fds[0].revents & POLLIN)) {
        char buf[128];
        ssize_t n = read(STDIN_FILENO, buf, sizeof buf - 1);
        if (n > 0) { buf[n] = '\0'; printf("Girdi: %s", buf); }
    } else if (hazir == 0) {
        printf("Zaman asimi.\n");
    }
    return 0;
}
\end{lstlisting}

\gsub{epoll --- Ölçeklenebilir Linux Çözümü}

\code{epoll}, on binlerce bağlantıyı verimli izler: her çağrıda tüm listeyi
taramaz, yalnızca \emph{hazır olanları} döndürür. Yüksek performanslı Linux
sunucularının (nginx, redis) temelidir.

\begin{lstlisting}
#include <sys/epoll.h>
#include <unistd.h>
#include <stdio.h>
int main(void) {
    int ep = epoll_create1(0);           // epoll ornegi olustur

    struct epoll_event olay = {0};
    olay.events  = EPOLLIN;              // okuma olaylarini izle
    olay.data.fd = STDIN_FILENO;
    epoll_ctl(ep, EPOLL_CTL_ADD, STDIN_FILENO, &olay);  // fd ekle

    struct epoll_event olaylar[10];
    printf("Giris bekleniyor...\n");

    // -1: sonsuza kadar bekle. Doner: hazir olay sayisi
    int n = epoll_wait(ep, olaylar, 10, -1);
    for (int i = 0; i < n; i++) {
        if (olaylar[i].data.fd == STDIN_FILENO) {
            char buf[128];
            ssize_t r = read(STDIN_FILENO, buf, sizeof buf - 1);
            if (r > 0) { buf[r] = '\0'; printf("Girdi: %s", buf); }
        }
    }
    close(ep);
    return 0;
}
\end{lstlisting}

\begin{center}
\begin{tabularx}{\textwidth}{@{}>{\bfseries\color{anaMavi}}l l >{\raggedright\arraybackslash}X@{}}
\toprule
\rowcolor{acikMavi}\color{anaMavi}\textbf{API} & \color{anaMavi}\textbf{Ölçek} & \color{anaMavi}\textbf{Not}\\
\midrule
select & Düşük (< 1024) & Taşınabilir (POSIX); küme her çağrıda kurulur.\\
\rowcolor{acikGri} poll & Orta & Taşınabilir; fd limiti yok ama O(n) tarama.\\
epoll & Yüksek (10K+) & Linux'a özgü; O(1), olay tetiklemeli.\\
\bottomrule
\end{tabularx}
\end{center}

\ipucu{\code{epoll} iki modda çalışır: \textbf{level-triggered} (varsayılan;
veri durdukça uyarır) ve \textbf{edge-triggered} (\code{EPOLLET}; yalnızca
durum \emph{değişince} uyarır). Edge-triggered daha hızlıdır ama tüm veriyi
non-blocking döngüde okumayı gerektirir.}

% =====================================================================
\gsec{Süreçler: fork, exec, wait}
% =====================================================================

Süreç (process), çalışan bir programdır. Linux'ta her süreç \code{fork} ile
oluşturulur ve \code{exec} ile başka bir programa dönüşebilir.

\gsub{fork: Süreç Çoğaltma}

\code{fork} çağıran süreci ikiye böler: ebeveyn (parent) ve neredeyse özdeş
bir çocuk (child). Tek çağrı, \emph{iki kez} döner.

\begin{lstlisting}
#include <unistd.h>
#include <stdio.h>
#include <sys/wait.h>
int main(void) {
    printf("fork oncesi (tek surec)\n");

    pid_t pid = fork();       // burada surec ikiye bolunur

    if (pid < 0) {            // hata
        perror("fork");
        return 1;
    } else if (pid == 0) {    // COCUK: fork 0 dondurur
        printf("Cocuk : PID=%d, ebeveyn=%d\n", getpid(), getppid());
    } else {                  // EBEVEYN: fork cocugun PID'sini dondurur
        printf("Ebeveyn: PID=%d, cocuk=%d\n", getpid(), pid);
        wait(NULL);           // cocugun bitmesini bekle (zombi olmasin)
    }
    return 0;
}
\end{lstlisting}

\begin{center}
\begin{tabularx}{\textwidth}{@{}>{\ttfamily\color{anaMavi}}l >{\raggedright\arraybackslash}X@{}}
\toprule
\rowcolor{acikMavi}\color{anaMavi}\textbf{fork dönüşü} & \color{anaMavi}\textbf{Anlam}\\
\midrule
< 0 & Hata (fork başarısız).\\
\rowcolor{acikGri} == 0 & Çocuk sürecindeyiz.\\
> 0 & Ebeveyndeyiz; değer çocuğun PID'sidir.\\
\bottomrule
\end{tabularx}
\end{center}

\bilgi{Çocuk, ebeveynin belleğinin bir \emph{kopyasını} alır (copy-on-write:
gerçek kopya ancak biri yazınca yapılır). Açık dosya betimleyicileri de
paylaşılır.}

\gsub{exec: Program Değiştirme}

\code{exec} ailesi, mevcut sürecin bellek imajını \emph{yeni bir programla}
tümüyle değiştirir. Başarılıysa geri dönmez.

\begin{center}
\begin{tabularx}{\textwidth}{@{}>{\ttfamily\color{anaMavi}}l >{\raggedright\arraybackslash}X@{}}
\toprule
\rowcolor{acikMavi}\color{anaMavi}\textbf{Fonksiyon} & \color{anaMavi}\textbf{Özellik (l=liste, v=vektör, p=PATH, e=env)}\\
\midrule
execl & Argümanlar liste olarak, tam yol.\\
\rowcolor{acikGri} execlp & Liste + PATH'te arar.\\
execv & Argümanlar dizi (vektör) olarak.\\
\rowcolor{acikGri} execvp & Vektör + PATH'te arar (en sık kullanılan).\\
execve & Vektör + özel ortam; \emph{gerçek} sistem çağrısı.\\
\bottomrule
\end{tabularx}
\end{center}

\ornek[fork + exec: kabuğun temel döngüsü]{Bir komut çalıştırmanın klasik
kalıbı: fork ile çocuk üret, çocukta exec ile komuta dönüş, ebeveynde bekle.}
\begin{lstlisting}
#include <unistd.h>
#include <sys/wait.h>
#include <stdio.h>
int main(void) {
    pid_t pid = fork();
    if (pid == 0) {
        // COCUK: "ls -l" komutuna donus
        char *argv[] = {"ls", "-l", NULL};   // NULL ile bitmeli
        execvp("ls", argv);
        // execvp basariliysa BURAYA ASLA GELINMEZ
        perror("execvp");                    // buraya gelinirse hata var
        _exit(127);                          // exec basarisiz
    } else if (pid > 0) {
        int durum;
        waitpid(pid, &durum, 0);             // cocugu bekle
        if (WIFEXITED(durum))
            printf("Cocuk %d cikti, kod=%d\n", pid, WEXITSTATUS(durum));
    }
    return 0;
}
\end{lstlisting}

\gsub{wait, waitpid ve Çıkış Durumu}

Ebeveyn, çocuğun çıkış durumunu \code{wait}/\code{waitpid} ile toplamalıdır.
Aksi halde çocuk "zombi" olarak süreç tablosunda kalır.

\begin{center}
\begin{tabularx}{\textwidth}{@{}>{\ttfamily\color{anaMavi}}l >{\raggedright\arraybackslash}X@{}}
\toprule
\rowcolor{acikMavi}\color{anaMavi}\textbf{Makro} & \color{anaMavi}\textbf{Anlam}\\
\midrule
WIFEXITED(s) & Normal mi çıktı (return/exit)?\\
\rowcolor{acikGri} WEXITSTATUS(s) & Çıkış kodu (0--255).\\
WIFSIGNALED(s) & Sinyalle mi öldü?\\
\rowcolor{acikGri} WTERMSIG(s) & Öldüren sinyal numarası.\\
\bottomrule
\end{tabularx}
\end{center}

\uyari{\textbf{Zombi süreç:} çocuk bitmiş ama ebeveyn \code{wait} etmemiştir;
çıkış durumu için tabloda yer tutar. \textbf{Öksüz (orphan) süreç:} ebeveyn
önce ölmüştür; çocuk \code{init}/\code{systemd} (PID 1) tarafından evlat
edinilir. Zombileri önlemek için her fork'a bir wait düşür ya da
\code{SIGCHLD} handler'ı kur.}

% =====================================================================
\gsec{İleri Süreç Yönetimi}
% =====================================================================

fork/exec/wait temellerinin ötesinde: exec ailesinin tümü, ortam değişkenleri,
çıkış kancaları, süreç grupları/oturumlar, arka plan servisleri (daemon) ve
kaynak sınırları.

\gsub{exec Ailesinin Tümü}

Altı \code{exec} varyantı, argümanları nasıl aldıklarına ve PATH/ortam
kullanıp kullanmadıklarına göre farklılaşır. Sonek harfleri: \textbf{l}=liste,
\textbf{v}=vektör, \textbf{p}=PATH'te ara, \textbf{e}=özel ortam.

\begin{lstlisting}
#include <unistd.h>
int main(void) {
    // l: argumanlar tek tek, NULL ile biter, TAM YOL
    execl("/bin/ls", "ls", "-l", "/tmp", (char *)NULL);

    // v: argumanlar dizi olarak
    char *av[] = {"ls", "-l", NULL};
    execv("/bin/ls", av);

    // p: komutu PATH'te arar (tam yol gerekmez) -- en pratik
    execlp("ls", "ls", "-l", (char *)NULL);
    execvp("ls", av);

    // e: cagirana ozel ortam (environment) verir
    char *env[] = {"YOL=/ozel", "DIL=tr", NULL};
    execle("/bin/prog", "prog", (char *)NULL, env);
    execvpe("prog", av, env);   // GNU eklentisi

    return 0;   // exec basarili olursa buraya asla gelinmez
}
\end{lstlisting}

\notk[Argüman listesi ve program adı]{exec'e verilen ilk argüman, çalışacak
programın \code{argv[0]}'ı olur --- genelde program adıdır ama farklı da
olabilir. Liste (\code{l}) biçiminde son argüman mutlaka \code{(char *)NULL}
olmalıdır.}

\gsub{Ortam Değişkenleri}

\begin{lstlisting}
#include <stdlib.h>
#include <stdio.h>
extern char **environ;         // tum ortam: "AD=deger" dizisi (NULL biter)

int main(void) {
    printf("HOME = %s\n", getenv("HOME"));   // oku

    setenv("UYGULAMA_MODU", "gelistirme", 1); // yaz (1 = ustune yaz)
    printf("MOD = %s\n", getenv("UYGULAMA_MODU"));

    unsetenv("UYGULAMA_MODU");                // sil

    // Tum ortami dolas
    for (char **e = environ; *e; e++)
        printf("%s\n", *e);
    return 0;
}
\end{lstlisting}

\gsub{Çıkış Kancaları: atexit ve \_exit}

\begin{lstlisting}
#include <stdlib.h>
#include <stdio.h>
#include <unistd.h>

void temizle1(void) { printf("temizlik 1\n"); }
void temizle2(void) { printf("temizlik 2\n"); }

int main(void) {
    atexit(temizle1);       // exit() ya da main return'de calisir
    atexit(temizle2);       // TERS sirada: once temizle2, sonra temizle1

    printf("is bitti\n");
    return 0;               // exit(0) cagrilir -> kancalar tetiklenir
    // _exit(0);            // FARK: kancalari CALISTIRMAZ, tampon bosaltmaz
}
\end{lstlisting}

\uyari{fork edilen çocukta hata olursa, \code{exit} yerine \code{\_exit}
(ya da \code{\_Exit}) kullan. \code{exit}, ebeveynden miras kalan stdio
tamponlarını boşaltarak çıktının \emph{iki kez} yazılmasına yol açabilir.}

\gsub{Süreç Grupları ve Oturumlar}

Her süreç bir \emph{süreç grubuna}, her grup bir \emph{oturuma} (session)
aittir. Kabuk, iş kontrolü (job control) ve sinyal dağıtımı bunlarla çalışır;
bir terminale bağlı tüm süreçlere aynı anda sinyal (örn. Ctrl+C) gönderilir.

\begin{lstlisting}
#include <unistd.h>
#include <stdio.h>
int main(void) {
    printf("PID=%d  PGID=%d  SID=%d\n",
           getpid(), getpgrp(), getsid(0));

    setpgid(0, 0);      // kendini yeni bir surec grubu lideri yap
    return 0;
}
\end{lstlisting}

\gsub{Arka Plan Servisi (Daemon) Oluşturma}

Bir daemon, terminalden kopmuş, arka planda çalışan bir servistir. Klasik
"double-fork" tekniğiyle oluşturulur.

\begin{lstlisting}
#include <unistd.h>
#include <stdlib.h>
#include <fcntl.h>
#include <sys/stat.h>

void daemon_ol(void) {
    // 1) Ilk fork: ebeveyn cikar, cocuk devam eder (kabuk komut istemine doner)
    if (fork() > 0) exit(0);

    // 2) Yeni oturum lideri ol: kontrol terminalinden kopar
    setsid();

    // 3) Ikinci fork: terminal edinmeyi kesin olarak engelle
    if (fork() > 0) exit(0);

    umask(0);               // dosya izin maskesini sifirla
    chdir("/");             // koke gec (bagli dizin unmount edilebilsin)

    // 4) Standart betimleyicileri /dev/null'a yonlendir
    int null = open("/dev/null", O_RDWR);
    dup2(null, STDIN_FILENO);
    dup2(null, STDOUT_FILENO);
    dup2(null, STDERR_FILENO);
    if (null > 2) close(null);
    // Artik arka planda calisan bir servisiz.
}
\end{lstlisting}

\ipucu{Modern sistemlerde daemon'ları elle yazmak yerine \code{systemd} servis
birimleri tercih edilir; systemd fork/setsid/log yönlendirmeyi kendi halleder.
Yine de tekniği bilmek altyapıyı anlamak için değerlidir.}

\gsub{Kaynak Sınırları: getrlimit / setrlimit}

\begin{lstlisting}
#include <sys/resource.h>
#include <stdio.h>
int main(void) {
    struct rlimit lim;
    getrlimit(RLIMIT_NOFILE, &lim);       // acik dosya siniri
    printf("Acik dosya: yumusak=%ld sert=%ld\n",
           (long)lim.rlim_cur, (long)lim.rlim_max);

    lim.rlim_cur = 256;                   // yumusak siniri dusur
    setrlimit(RLIMIT_NOFILE, &lim);

    // RLIMIT_CPU (cpu saniyesi), RLIMIT_AS (adres alani),
    // RLIMIT_CORE (core dump boyutu) de ayarlanabilir.
    return 0;
}
\end{lstlisting}

% =====================================================================
\gsec{Terminal Komutları ve Program Çalıştırma}
% =====================================================================

Bir C programından kabuk komutu ya da başka bir program çalıştırmanın çeşitli
yolları: en basitten (\code{system}) en güçlüye (elle fork/exec + boru ile
çıktı yakalama) kadar.

\gsub{system: En Basit Yol (ve Tehlikeleri)}

\code{system}, komutu \code{/bin/sh -c} ile çalıştırır ve bitmesini bekler.
Basittir ama kabuk enjeksiyonuna (shell injection) açıktır.

\begin{lstlisting}
#include <stdlib.h>
#include <stdio.h>
#include <sys/wait.h>
int main(void) {
    int durum = system("ls -l /tmp | wc -l");   // kabuk ozelliklerini kullanir

    if (durum == -1) perror("system");
    else printf("Cikis kodu: %d\n", WEXITSTATUS(durum));

    return 0;
}
\end{lstlisting}

\uyari{Kullanıcı girdisini asla doğrudan \code{system}'e verme:
\code{system(strcat("rm ", kullanici\_girdisi))} felakettir --- kullanıcı
\code{"; rm -rf /"} yazarsa çalışır. Güvenlik gerektiğinde \code{fork} +
\code{execvp} kullan (kabuk araya girmez, argümanlar yorumlanmaz).}

\gsub{popen: Komut Çıktısını Okuma}

\code{popen}, komutu çalıştırıp çıktısını (ya da girdisini) bir \code{FILE*}
akışı olarak verir --- komutun sonucunu satır satır okumak için pratiktir.

\begin{lstlisting}
#include <stdio.h>
int main(void) {
    // "r": komutun STDOUT'unu oku. ("w": komutun STDIN'ine yaz)
    FILE *boru = popen("ls -1 /etc | head -5", "r");
    if (!boru) { perror("popen"); return 1; }

    char satir[256];
    while (fgets(satir, sizeof satir, boru))   // ciktiyi satir satir oku
        printf("--> %s", satir);

    int durum = pclose(boru);                  // kapat + cikis kodunu al
    printf("Komut cikis kodu: %d\n", WEXITSTATUS(durum));
    return 0;
}
\end{lstlisting}

\gsub{Elle Çıktı Yakalama: fork + exec + pipe}

\code{popen} kabuk kullanır; tam kontrol ve güvenlik için boruyu kendin kurarsın.
Bu, bir komutu çalıştırıp çıktısını tampona toplamanın "endüstriyel" yoludur.

\begin{lstlisting}
#include <unistd.h>
#include <sys/wait.h>
#include <stdio.h>
#include <string.h>

int main(void) {
    int boru[2];
    if (pipe(boru) == -1) { perror("pipe"); return 1; }

    pid_t pid = fork();
    if (pid == 0) {                       // COCUK: komutu calistir
        close(boru[0]);                   // okuma ucu gereksiz
        dup2(boru[1], STDOUT_FILENO);     // stdout -> borunun yazma ucu
        close(boru[1]);
        char *av[] = {"date", "+%Y-%m-%d %H:%M:%S", NULL};
        execvp("date", av);               // kabuk YOK: guvenli
        _exit(127);                       // exec basarisizsa
    }

    // EBEVEYN: cocugun ciktisini oku
    close(boru[1]);                       // yazma ucu gereksiz
    char cikti[512];
    ssize_t toplam = 0, n;
    while ((n = read(boru[0], cikti + toplam, sizeof cikti - 1 - toplam)) > 0)
        toplam += n;
    cikti[toplam] = '\0';
    close(boru[0]);

    int durum;
    waitpid(pid, &durum, 0);              // zombi birakma
    printf("Yakalanan cikti: %s", cikti);
    printf("Cikis kodu: %d\n", WEXITSTATUS(durum));
    return 0;
}
\end{lstlisting}

\bilgi{Bu desen \code{system}/\code{popen}'a göre daha uzundur ama: (1) kabuk
enjeksiyonu riski yoktur, (2) çıkış kodunu ve stderr'i ayrı yakalayabilirsin,
(3) zaman aşımı/sinyal kontrolü sende olur. Üretim kodunda tercih edilen yol
budur.}

\gsub{Basit Bir Kabuk (Mini-Shell)}

Öğrenilen her şeyi birleştiren klasik alıştırma: komut oku, ayrıştır,
fork/exec ile çalıştır, bekle --- döngü.

\begin{lstlisting}
#include <stdio.h>
#include <stdlib.h>
#include <string.h>
#include <unistd.h>
#include <sys/wait.h>

int main(void) {
    char satir[1024];
    while (1) {
        printf("minish$ ");
        fflush(stdout);
        if (!fgets(satir, sizeof satir, stdin)) break;  // Ctrl+D -> cik
        satir[strcspn(satir, "\n")] = '\0';             // sondaki \n'i at
        if (satir[0] == '\0') continue;
        if (strcmp(satir, "exit") == 0) break;          // dahili komut

        // Bosluklara gore argumanlara ayir
        char *argv[64]; int argc = 0;
        char *tok = strtok(satir, " ");
        while (tok && argc < 63) { argv[argc++] = tok; tok = strtok(NULL, " "); }
        argv[argc] = NULL;

        pid_t pid = fork();
        if (pid == 0) {                    // COCUK
            execvp(argv[0], argv);
            perror(argv[0]);               // komut bulunamadi
            _exit(127);
        } else if (pid > 0) {              // EBEVEYN
            int durum;
            waitpid(pid, &durum, 0);       // komutun bitmesini bekle
        }
    }
    printf("\nHosca kal.\n");
    return 0;
}
\end{lstlisting}

\ipucu{Bu mini-shell'i geliştirmek harika bir egzersizdir: boru (\code{|}),
yönlendirme (\code{>}, \code{<}, \code{2>}), arka plan çalıştırma (\code{\&}),
\code{cd} gibi dahili komutlar ve sinyal yönetimi ekleyerek gerçek bir kabuğa
dönüştürebilirsin.}

% =====================================================================
\gsec{Sinyaller}
% =====================================================================

Sinyal, bir sürece gönderilen asenkron bildirimdir: Ctrl+C (\code{SIGINT}),
zamanlayıcı (\code{SIGALRM}), segfault (\code{SIGSEGV}), çocuk bitişi
(\code{SIGCHLD}) gibi.

\gsub{Yaygın Sinyaller}

\begin{center}
\begin{tabularx}{\textwidth}{@{}>{\ttfamily\color{anaMavi}}l l >{\raggedright\arraybackslash}X@{}}
\toprule
\rowcolor{acikMavi}\color{anaMavi}\textbf{Sinyal} & \color{anaMavi}\textbf{No} & \color{anaMavi}\textbf{Anlam / Varsayılan}\\
\midrule
SIGINT & 2 & Ctrl+C; sonlandır.\\
\rowcolor{acikGri} SIGKILL & 9 & Zorla öldür; \emph{yakalanamaz/engellenemez}.\\
SIGSEGV & 11 & Geçersiz bellek erişimi; core dump.\\
\rowcolor{acikGri} SIGTERM & 15 & Kibar sonlandırma isteği (varsayılan kill).\\
SIGCHLD & 17 & Çocuk süreç durdu/bitti.\\
\rowcolor{acikGri} SIGALRM & 14 & alarm() zamanlayıcısı doldu.\\
SIGSTOP / SIGCONT & 19/18 & Durdur / devam ettir.\\
\bottomrule
\end{tabularx}
\end{center}

\gsub{sigaction ile Sinyal Yakalama}

\code{signal()} taşınabilirlik açısından güvenilmezdir; modern kod
\code{sigaction} kullanır.

\begin{lstlisting}
#include <stdio.h>
#include <signal.h>
#include <unistd.h>

volatile sig_atomic_t calisiyor = 1;   // handler ile paylasilan bayrak

void handler(int sinyal) {
    // Handler icinde yalnizca async-signal-safe islemler yapilmali.
    // Bir bayrak set etmek guvenlidir; printf DEGILDIR.
    calisiyor = 0;
    (void)sinyal;
}

int main(void) {
    struct sigaction sa = {0};
    sa.sa_handler = handler;
    sigemptyset(&sa.sa_mask);          // handler sirasinda bloklanacak sinyaller
    sa.sa_flags = SA_RESTART;          // kesilen sistem cagrilarini otomatik tekrarla
    sigaction(SIGINT, &sa, NULL);      // Ctrl+C'yi yakala

    printf("Calisiyor... Ctrl+C ile durdur.\n");
    while (calisiyor) {
        pause();                       // bir sinyal gelene kadar uyu
    }
    printf("\nDuzgun sekilde cikildi.\n");
    return 0;
}
\end{lstlisting}

\uyari{Sinyal handler'ı içinde \emph{yalnızca} async-signal-safe fonksiyonlar
çağrılabilir (\code{write} evet, \code{printf}/\code{malloc} \textbf{hayır}).
Handler ile ana kod arasında paylaşılan değişken mutlaka
\code{volatile sig\_atomic\_t} olmalıdır.}

\gsub{kill ve alarm}

\begin{lstlisting}
#include <signal.h>
#include <unistd.h>
#include <stdio.h>

void zaman_asimi(int s) { (void)s; write(2, "Zaman doldu!\n", 13); }

int main(void) {
    struct sigaction sa = {0};
    sa.sa_handler = zaman_asimi;
    sigaction(SIGALRM, &sa, NULL);

    alarm(3);            // 3 saniye sonra SIGALRM gonder
    printf("Bir sey bekleniyor (en fazla 3 sn)...\n");
    pause();             // sinyali bekle

    // Baska bir surece sinyal gondermek:
    // kill(hedef_pid, SIGTERM);
    return 0;
}
\end{lstlisting}

\gsub{SIGCHLD ile Zombi Önleme}

\begin{lstlisting}
#include <signal.h>
#include <sys/wait.h>

void cocuk_bitti(int s) {
    (void)s;
    // Biriken tum bitmis cocuklari topla (WNOHANG: bloklanma)
    while (waitpid(-1, NULL, WNOHANG) > 0)
        ;
}
// main icinde: sigaction ile SIGCHLD -> cocuk_bitti kurulur.
// Boylece cocuklar otomatik toplanir, zombi kalmaz.
\end{lstlisting}

% =====================================================================
\gsec{Borular (Pipes) ve FIFO}
% =====================================================================

Boru, tek yönlü bir bayt akışıdır: bir uçtan yazılır, diğerinden okunur.
İlişkili süreçler (fork ile) arası iletişimin en basit yoludur.

\gsub{Anonim Boru (pipe)}

\begin{lstlisting}
#include <unistd.h>
#include <stdio.h>
#include <string.h>
int main(void) {
    int fd[2];               // fd[0]=okuma ucu, fd[1]=yazma ucu
    if (pipe(fd) == -1) { perror("pipe"); return 1; }

    pid_t pid = fork();
    if (pid == 0) {                     // COCUK: yazar
        close(fd[0]);                   // okuma ucunu kapat
        const char *mesaj = "Merhaba ebeveyn!";
        write(fd[1], mesaj, strlen(mesaj));
        close(fd[1]);
    } else {                            // EBEVEYN: okur
        close(fd[1]);                   // yazma ucunu kapat
        char tampon[128];
        ssize_t n = read(fd[0], tampon, sizeof(tampon) - 1);
        if (n > 0) { tampon[n] = '\0'; printf("Alindi: %s\n", tampon); }
        close(fd[0]);
        wait(NULL);
    }
    return 0;
}
\end{lstlisting}

\uyari{Kullanmadığın boru uçlarını \emph{mutlaka kapat}. Yazma ucu tümüyle
kapanmazsa okuyan taraf EOF (0) alamaz ve sonsuza kadar bekler. Bu, en sık
yapılan boru hatasıdır.}

\gsub{Boru + exec: Kabuğun "\code{|}" Operatörü}

\code{dup2} ile bir sürecin stdout'unu diğerinin stdin'ine bağlayarak
\code{ls | wc -l} eşdeğerini kurabiliriz.

\begin{lstlisting}
#include <unistd.h>
#include <sys/wait.h>
int main(void) {
    int fd[2];
    pipe(fd);

    if (fork() == 0) {                  // 1. cocuk: "ls"
        dup2(fd[1], STDOUT_FILENO);     // stdout -> borunun yazma ucu
        close(fd[0]); close(fd[1]);
        execlp("ls", "ls", NULL);
        _exit(127);
    }
    if (fork() == 0) {                  // 2. cocuk: "wc -l"
        dup2(fd[0], STDIN_FILENO);      // stdin -> borunun okuma ucu
        close(fd[0]); close(fd[1]);
        execlp("wc", "wc", "-l", NULL);
        _exit(127);
    }
    close(fd[0]); close(fd[1]);         // ebeveyn iki ucu da kapatir
    wait(NULL); wait(NULL);             // iki cocugu bekle
    return 0;
}
\end{lstlisting}

\gsub{Adlandırılmış Boru (FIFO)}

FIFO, dosya sisteminde bir isme sahip boru olduğu için \emph{ilişkisiz}
süreçler de kullanabilir.

\begin{lstlisting}
// --- yazan.c ---
#include <sys/stat.h>
#include <fcntl.h>
#include <unistd.h>
int main(void) {
    mkfifo("/tmp/borum", 0666);          // FIFO olustur (bir kez yeter)
    int fd = open("/tmp/borum", O_WRONLY); // okuyucu baglanana kadar bloke
    write(fd, "FIFO uzerinden selam", 20);
    close(fd);
    return 0;
}
// --- okuyan.c ---
// int fd = open("/tmp/borum", O_RDONLY);
// char buf[64]; ssize_t n = read(fd, buf, sizeof buf); ...
\end{lstlisting}

\gsub{Çift Yönlü İletişim (İki Boru)}

Bir boru tek yönlüdür. İki süreç arasında karşılıklı konuşma için \emph{iki}
boru kurulur --- biri her yön için.

\begin{lstlisting}
#include <unistd.h>
#include <sys/wait.h>
#include <stdio.h>
#include <string.h>

int main(void) {
    int eb_co[2], co_eb[2];       // ebeveyn->cocuk, cocuk->ebeveyn
    pipe(eb_co); pipe(co_eb);

    if (fork() == 0) {                       // COCUK
        close(eb_co[1]); close(co_eb[0]);    // kullanilmayan uclari kapat
        char buf[64];
        ssize_t n = read(eb_co[0], buf, sizeof buf - 1);
        buf[n] = '\0';
        printf("Cocuk aldi: %s\n", buf);
        write(co_eb[1], "selam ebeveyn", 13);  // cevap yaz
        _exit(0);
    } else {                                 // EBEVEYN
        close(eb_co[0]); close(co_eb[1]);
        write(eb_co[1], "selam cocuk", 11);  // once yaz
        char buf[64];
        ssize_t n = read(co_eb[0], buf, sizeof buf - 1);
        buf[n] = '\0';
        printf("Ebeveyn aldi: %s\n", buf);
        wait(NULL);
    }
    return 0;
}
\end{lstlisting}

\uyari{Çift yönlü boruda \textbf{kilitlenme} (deadlock) riski vardır: iki taraf
da aynı anda okuma beklerse ya da tampon dolarken ikisi de yazmaya çalışırsa
program donar. Protokolü net tanımla (kim önce yazar/okur) ve gerekirse
non-blocking G/Ç kullan.}

\gsub{N Komutluk Boru Hattı}

Kabuğun \code{a | b | c} zincirini kurmak için N-1 boru gerekir; her komut
öncekinin çıktısından okuyup sonrakine yazar.

\begin{lstlisting}
#include <unistd.h>
#include <sys/wait.h>
#include <stdio.h>

// Ornek: cat /etc/passwd | grep root | wc -l
int main(void) {
    char *komutlar[][4] = {
        {"cat", "/etc/passwd", NULL},
        {"grep", "root", NULL},
        {"wc", "-l", NULL},
    };
    int adet = 3;
    int onceki_okuma = -1;                 // ilk komutun girisi normal stdin

    for (int i = 0; i < adet; i++) {
        int boru[2];
        int son = (i == adet - 1);
        if (!son) pipe(boru);              // son komut disinda boru ac

        if (fork() == 0) {                 // COCUK
            if (onceki_okuma != -1) {      // ortadaki/son: onceki borudan oku
                dup2(onceki_okuma, STDIN_FILENO);
                close(onceki_okuma);
            }
            if (!son) {                    // son degilse: bir sonraki boruya yaz
                close(boru[0]);
                dup2(boru[1], STDOUT_FILENO);
                close(boru[1]);
            }
            execvp(komutlar[i][0], komutlar[i]);
            _exit(127);
        }
        // EBEVEYN: kullanilmayan uclari kapat, sonraki iterasyona hazirla
        if (onceki_okuma != -1) close(onceki_okuma);
        if (!son) { close(boru[1]); onceki_okuma = boru[0]; }
    }
    for (int i = 0; i < adet; i++) wait(NULL);  // tum cocuklari bekle
    return 0;
}
\end{lstlisting}

\gsub{SIGPIPE: Kapalı Boruya Yazma}

Okuma ucu kapalı bir boruya yazmak \code{SIGPIPE} sinyali üretir; varsayılan
davranış programı \emph{öldürmektir}. Ağ ve boru kodunda bu sinyal yönetilmelidir.

\begin{lstlisting}
#include <signal.h>
#include <unistd.h>
#include <errno.h>
#include <stdio.h>
int main(void) {
    signal(SIGPIPE, SIG_IGN);    // sinyali yok say -> write EPIPE dondursun

    int boru[2];
    pipe(boru);
    close(boru[0]);              // okuma ucunu kapat

    ssize_t r = write(boru[1], "veri", 4);   // kimse okumuyor
    if (r == -1 && errno == EPIPE)
        printf("Yazma basarisiz: karsi uc kapali (EPIPE)\n");
    close(boru[1]);
    return 0;
}
\end{lstlisting}

\notk[PIPE\_BUF ve Atomiklik]{\code{PIPE\_BUF} baytından (Linux'ta 4096) küçük
yazmalar \emph{atomiktir}: birden çok yazıcı olsa bile veriler iç içe geçmez.
Daha büyük yazmalar bölünebilir. Çok yazıcılı boru kayıtlarını bu sınırın
altında tut.}

\gsub{popen ile Boru --- Hızlı Alternatif}

Süreçler arası boruyu elle kurmak yerine, tek bir kabuk komutu için
\code{popen} çok daha kısadır (Terminal Komutları bölümüne bakınız). Elle
\code{pipe}+\code{dup2} ise kabuk kullanmadan tam kontrol sağlar --- güvenlik
ve esneklik gerektiğinde onu seç.

% =====================================================================
\gsec{IPC: Paylaşımlı Bellek, Mesaj Kuyruğu, Semafor}
% =====================================================================

Süreçler arası iletişim (Inter-Process Communication) için borular dışında
üç güçlü mekanizma vardır. Modern kod POSIX API'lerini tercih eder.

\gsub{POSIX Paylaşımlı Bellek (shm\_open + mmap)}

En hızlı IPC'dir: iki süreç aynı fiziksel belleği görür; kopyalama yoktur.

\begin{lstlisting}
#include <sys/mman.h>
#include <sys/stat.h>
#include <fcntl.h>
#include <unistd.h>
#include <stdio.h>
#include <string.h>

#define AD "/paylasim"
#define BOYUT 4096

int main(void) {
    // Paylasimli bellek nesnesi olustur/ac
    int fd = shm_open(AD, O_CREAT | O_RDWR, 0666);
    ftruncate(fd, BOYUT);                       // boyutlandir

    // Belegi surecin adres alanina esle
    char *bellek = mmap(NULL, BOYUT, PROT_READ | PROT_WRITE,
                        MAP_SHARED, fd, 0);
    if (bellek == MAP_FAILED) { perror("mmap"); return 1; }

    if (fork() == 0) {                          // COCUK: yazar
        strcpy(bellek, "Paylasimli bellekten selam");
        _exit(0);
    } else {                                    // EBEVEYN: okur
        wait(NULL);
        printf("Okundu: %s\n", bellek);
    }

    munmap(bellek, BOYUT);
    close(fd);
    shm_unlink(AD);                             // isimlendirilmis nesneyi sil
    return 0;
}
// Derleme: gcc bu.c -o bu -lrt
\end{lstlisting}

\uyari{Paylaşımlı bellekte \emph{senkronizasyon yoktur}: iki süreç aynı anda
yazarsa yarış durumu (race condition) oluşur. Erişimi bir semafor veya mutex
ile korumalısın.}

\gsub{POSIX Semafor}

Semafor, paylaşılan bir kaynağa erişimi sayan bir kilittir. \code{sem\_wait}
(azalt/bekle) ve \code{sem\_post} (artır/serbest bırak) atomiktir.

\begin{lstlisting}
#include <semaphore.h>
#include <fcntl.h>
#include <stdio.h>
#include <unistd.h>
#include <sys/wait.h>

int main(void) {
    // Baslangic degeri 1 -> ikili (mutex gibi) semafor
    sem_t *sem = sem_open("/semaforum", O_CREAT, 0666, 1);

    if (fork() == 0) {
        sem_wait(sem);                 // kritik bolgeye gir (kilitle)
        printf("Cocuk kritik bolgede\n");
        sleep(1);
        sem_post(sem);                 // cik (birak)
        _exit(0);
    } else {
        sem_wait(sem);
        printf("Ebeveyn kritik bolgede\n");
        sem_post(sem);
        wait(NULL);
    }
    sem_close(sem);
    sem_unlink("/semaforum");
    return 0;
}
// Derleme: gcc bu.c -o bu -lpthread -lrt
\end{lstlisting}

\gsub{POSIX Mesaj Kuyruğu}

Mesaj kuyruğu, süreçler arasında \emph{sınır korumalı} (mesaj mesaj) ve
öncelikli veri aktarımı sağlar.

\begin{lstlisting}
#include <mqueue.h>
#include <stdio.h>
#include <string.h>
int main(void) {
    struct mq_attr attr = {0};
    attr.mq_maxmsg = 10;               // kuyrukta en fazla 10 mesaj
    attr.mq_msgsize = 256;             // her mesaj en fazla 256 bayt

    mqd_t mq = mq_open("/kuyruk", O_CREAT | O_RDWR, 0666, &attr);

    mq_send(mq, "birinci", 7, 1);      // (kuyruk, veri, uzunluk, oncelik)
    mq_send(mq, "onemli", 6, 9);       // yuksek oncelik once alinir

    char buf[256]; unsigned oncelik;
    ssize_t n = mq_receive(mq, buf, sizeof buf, &oncelik);
    printf("Alindi (%u): %.*s\n", oncelik, (int)n, buf); // onemli

    mq_close(mq);
    mq_unlink("/kuyruk");
    return 0;
}
// Derleme: gcc bu.c -o bu -lrt
\end{lstlisting}

\begin{center}
\begin{tabularx}{\textwidth}{@{}>{\bfseries\color{anaMavi}}l >{\raggedright\arraybackslash}X >{\raggedright\arraybackslash}X@{}}
\toprule
\rowcolor{acikMavi}\color{anaMavi}\textbf{Mekanizma} & \color{anaMavi}\textbf{Güçlü Yanı} & \color{anaMavi}\textbf{Not}\\
\midrule
Boru/FIFO & Basit, akış temelli & Tek yönlü, mesaj sınırı yok\\
\rowcolor{acikGri} Paylaşımlı bellek & En hızlı & Senkronizasyon gerekir\\
Mesaj kuyruğu & Mesaj sınırı + öncelik & Kopyalama maliyeti\\
\rowcolor{acikGri} Semafor & Senkronizasyon & Veri taşımaz, koordine eder\\
\bottomrule
\end{tabularx}
\end{center}

% =====================================================================
\gsec{İş Parçacıkları (POSIX Threads)}
% =====================================================================

İş parçacığı (thread), aynı süreç içinde paralel çalışan bir yürütme
akışıdır. Süreçlerin aksine iş parçacıkları \emph{aynı bellek alanını}
paylaşır --- bu hem güçlü hem tehlikelidir.

\gsub{İş Parçacığı Oluşturma}

\begin{lstlisting}
#include <pthread.h>
#include <stdio.h>

void *isci(void *arg) {
    int id = *(int *)arg;
    printf("Isci %d calisiyor\n", id);
    return NULL;              // pthread_exit(NULL) ile ayni
}

int main(void) {
    pthread_t iplikler[3];
    int idler[3] = {0, 1, 2};

    for (int i = 0; i < 3; i++)
        pthread_create(&iplikler[i], NULL, isci, &idler[i]);

    for (int i = 0; i < 3; i++)
        pthread_join(iplikler[i], NULL);   // her ipligin bitmesini bekle

    printf("Tum iplikler bitti\n");
    return 0;
}
// Derleme: gcc bu.c -o bu -lpthread
\end{lstlisting}

\gsub{Yarış Durumu (Race Condition)}

Paylaşılan bir değişkene senkronizasyonsuz erişim, öngörülemez sonuç verir.

\begin{lstlisting}
#include <pthread.h>
#include <stdio.h>
long sayac = 0;                     // paylasilan; KORUNMASIZ

void *artir(void *arg) {
    (void)arg;
    for (int i = 0; i < 1000000; i++)
        sayac++;                    // sayac++ atomik DEGIL: oku-artir-yaz
    return NULL;
}
int main(void) {
    pthread_t a, b;
    pthread_create(&a, NULL, artir, NULL);
    pthread_create(&b, NULL, artir, NULL);
    pthread_join(a, NULL); pthread_join(b, NULL);
    printf("sayac = %ld (beklenen 2000000)\n", sayac); // genelde daha az!
    return 0;
}
\end{lstlisting}

\gsub{Mutex ile Koruma}

\begin{lstlisting}
#include <pthread.h>
#include <stdio.h>
long sayac = 0;
pthread_mutex_t kilit = PTHREAD_MUTEX_INITIALIZER;

void *artir(void *arg) {
    (void)arg;
    for (int i = 0; i < 1000000; i++) {
        pthread_mutex_lock(&kilit);      // kritik bolgeyi kilitle
        sayac++;
        pthread_mutex_unlock(&kilit);    // birak
    }
    return NULL;
}
int main(void) {
    pthread_t a, b;
    pthread_create(&a, NULL, artir, NULL);
    pthread_create(&b, NULL, artir, NULL);
    pthread_join(a, NULL); pthread_join(b, NULL);
    printf("sayac = %ld\n", sayac);      // her zaman 2000000
    pthread_mutex_destroy(&kilit);
    return 0;
}
\end{lstlisting}

\uyari{\textbf{Deadlock (kilitlenme):} iki iş parçacığı birbirinin tuttuğu
kilidi beklerse ikisi de sonsuza kadar durur. Önlem: tüm kilitleri
\emph{her yerde aynı sırada} al ve kritik bölgeyi kısa tut.}

\gsub{Koşul Değişkeni (Condition Variable)}

Üretici-tüketici gibi senaryolarda, bir iş parçacığının bir koşulu beklemesi
için \code{pthread\_cond\_t} kullanılır.

\begin{lstlisting}
#include <pthread.h>
#include <stdio.h>

pthread_mutex_t m = PTHREAD_MUTEX_INITIALIZER;
pthread_cond_t  c = PTHREAD_COND_INITIALIZER;
int hazir = 0;

void *tuketici(void *arg) {
    (void)arg;
    pthread_mutex_lock(&m);
    while (!hazir)                       // DAIMA dongude bekle (spurious wakeup)
        pthread_cond_wait(&c, &m);       // kilidi birakip uyur, uyaninca kilitler
    printf("Tuketici: veri hazir!\n");
    pthread_mutex_unlock(&m);
    return NULL;
}
void *uretici(void *arg) {
    (void)arg;
    pthread_mutex_lock(&m);
    hazir = 1;
    pthread_cond_signal(&c);             // bekleyeni uyandir
    pthread_mutex_unlock(&m);
    return NULL;
}
int main(void) {
    pthread_t t, u;
    pthread_create(&t, NULL, tuketici, NULL);
    pthread_create(&u, NULL, uretici, NULL);
    pthread_join(t, NULL); pthread_join(u, NULL);
    return 0;
}
\end{lstlisting}

\ipucu{\code{pthread\_cond\_wait} her zaman bir \code{while} döngüsü içinde
kullanılır çünkü "sahte uyanma" (spurious wakeup) mümkündür. \code{if} ile
kontrol \emph{hatalıdır}.}

\gsub{İş Parçacığı mı, Süreç mi?}

\begin{center}
\begin{tabularx}{\textwidth}{@{}>{\raggedright\arraybackslash}X >{\raggedright\arraybackslash}X@{}}
\toprule
\rowcolor{acikMavi}\color{anaMavi}\textbf{İş Parçacığı (thread)} & \color{anaMavi}\textbf{Süreç (process)}\\
\midrule
Belleği paylaşır (hızlı iletişim) & Ayrı bellek (izolasyon, güvenlik)\\
\rowcolor{acikGri} Oluşturması ucuz & Oluşturması pahalı (fork)\\
Biri çökerse tümü çöker & Biri çökerse diğerleri sağ kalır\\
\rowcolor{acikGri} Senkronizasyon şart & Doğal olarak yalıtılmış\\
\bottomrule
\end{tabularx}
\end{center}

% =====================================================================
\gsec{Ağ Programlama: Soketler}
% =====================================================================

Soket, ağ üzerinden (ya da aynı makinede) iletişim uç noktasıdır. TCP
güvenilir bağlantı akışı, UDP ise bağlantısız datagram sunar.

\gsub{TCP Sunucu}

\begin{lstlisting}
#include <sys/socket.h>
#include <netinet/in.h>
#include <arpa/inet.h>
#include <unistd.h>
#include <string.h>
#include <stdio.h>

int main(void) {
    // 1) Soket olustur (IPv4, TCP)
    int s = socket(AF_INET, SOCK_STREAM, 0);
    if (s == -1) { perror("socket"); return 1; }

    // Adresi hemen tekrar kullanilabilir yap (TIME_WAIT sorununu asar)
    int opt = 1;
    setsockopt(s, SOL_SOCKET, SO_REUSEADDR, &opt, sizeof opt);

    // 2) Adrese bagla (bind)
    struct sockaddr_in adr = {0};
    adr.sin_family = AF_INET;
    adr.sin_addr.s_addr = INADDR_ANY;      // tum arayuzler
    adr.sin_port = htons(8080);            // port (ag bayt sirasi!)
    if (bind(s, (struct sockaddr *)&adr, sizeof adr) == -1) {
        perror("bind"); return 1;
    }

    // 3) Dinlemeye basla (listen)
    listen(s, 10);                         // 10 baglanti kuyruk boyu
    printf("8080 portu dinleniyor...\n");

    // 4) Baglanti kabul et (accept) -> yeni bir fd doner
    struct sockaddr_in istemci;
    socklen_t len = sizeof istemci;
    int c = accept(s, (struct sockaddr *)&istemci, &len);

    char ip[INET_ADDRSTRLEN];
    inet_ntop(AF_INET, &istemci.sin_addr, ip, sizeof ip);
    printf("Baglanti: %s:%d\n", ip, ntohs(istemci.sin_port));

    // 5) Veri al/gonder
    char tampon[1024];
    ssize_t n = recv(c, tampon, sizeof tampon - 1, 0);
    if (n > 0) { tampon[n] = '\0'; printf("Istemci: %s\n", tampon); }
    const char *cevap = "Merhaba istemci!\n";
    send(c, cevap, strlen(cevap), 0);

    close(c);
    close(s);
    return 0;
}
\end{lstlisting}

\gsub{TCP İstemci}

\begin{lstlisting}
#include <sys/socket.h>
#include <netinet/in.h>
#include <arpa/inet.h>
#include <unistd.h>
#include <string.h>
#include <stdio.h>

int main(void) {
    int s = socket(AF_INET, SOCK_STREAM, 0);

    struct sockaddr_in sunucu = {0};
    sunucu.sin_family = AF_INET;
    sunucu.sin_port = htons(8080);
    inet_pton(AF_INET, "127.0.0.1", &sunucu.sin_addr);  // metin -> ikili IP

    if (connect(s, (struct sockaddr *)&sunucu, sizeof sunucu) == -1) {
        perror("connect"); return 1;
    }

    const char *mesaj = "Sunucuya selam";
    send(s, mesaj, strlen(mesaj), 0);

    char tampon[1024];
    ssize_t n = recv(s, tampon, sizeof tampon - 1, 0);
    if (n > 0) { tampon[n] = '\0'; printf("Sunucu: %s", tampon); }

    close(s);
    return 0;
}
\end{lstlisting}

\begin{center}
\begin{tabularx}{\textwidth}{@{}>{\bfseries\color{anaMavi}}l >{\raggedright\arraybackslash}X@{}}
\toprule
\rowcolor{acikMavi}\color{anaMavi}\textbf{Sunucu Adımı} & \color{anaMavi}\textbf{İstemci Adımı}\\
\midrule
socket() $\to$ bind() $\to$ listen() $\to$ accept() & socket() $\to$ connect()\\
\bottomrule
\end{tabularx}
\end{center}

\uyari{Ağ üzerindeki tüm çok baytlı sayılar \emph{ağ bayt sırasında}
(big-endian) taşınır. Port ve adresleri \code{htons}/\code{htonl} ile gönder,
\code{ntohs}/\code{ntohl} ile oku. Aksi halde farklı mimariler arası iletişim
bozulur.}

\gsub{UDP: Bağlantısız İletişim}

\begin{lstlisting}
// --- UDP Sunucu (ozet) ---
int s = socket(AF_INET, SOCK_DGRAM, 0);   // SOCK_DGRAM = UDP
// bind(...) ayni sekilde
struct sockaddr_in istemci; socklen_t len = sizeof istemci;
char buf[1024];
ssize_t n = recvfrom(s, buf, sizeof buf, 0,
                     (struct sockaddr *)&istemci, &len);  // kimden geldi?
sendto(s, "cevap", 5, 0, (struct sockaddr *)&istemci, len); // ona geri gonder
// UDP'de listen/accept/connect YOK; her datagram bagimsizdir.
\end{lstlisting}

% =====================================================================
\gsec{mmap: Bellek Eşleme}
% =====================================================================

\code{mmap}, bir dosyayı ya da anonim bir bellek bölgesini doğrudan sürecin
adres alanına eşler. Dosyaya \code{read}/\code{write} yerine \emph{diziymiş
gibi} erişilir --- büyük dosyalar için çok verimlidir.

\begin{lstlisting}
#include <sys/mman.h>
#include <sys/stat.h>
#include <fcntl.h>
#include <unistd.h>
#include <stdio.h>

int main(int argc, char *argv[]) {
    if (argc != 2) { fprintf(stderr, "Kullanim: %s dosya\n", argv[0]); return 1; }

    int fd = open(argv[1], O_RDONLY);
    if (fd == -1) { perror("open"); return 1; }

    struct stat st;
    fstat(fd, &st);                          // dosya boyutunu al

    // Dosyayi salt-okunur olarak bellege esle
    char *harita = mmap(NULL, st.st_size, PROT_READ, MAP_PRIVATE, fd, 0);
    if (harita == MAP_FAILED) { perror("mmap"); close(fd); return 1; }
    close(fd);                               // esleme sonrasi fd gereksiz

    // Artik dosya bir char dizisi gibi: satirlari say
    long satir = 0;
    for (off_t i = 0; i < st.st_size; i++)
        if (harita[i] == '\n') satir++;
    printf("Satir sayisi: %ld\n", satir);

    munmap(harita, st.st_size);              // eslemeyi kaldir
    return 0;
}
\end{lstlisting}

\bilgi{\code{mmap} ayrıca \code{MAP\_ANONYMOUS} ile dosyasız bellek ayırmak
(malloc'un altında kullanılır) ve \code{MAP\_SHARED} ile süreçler arası
paylaşımlı bellek için de kullanılır. \code{PROT\_WRITE} ile eşlenirse
belleğe yazmak dosyayı değiştirir.}

% =====================================================================
\gsec{Zaman ve Zamanlayıcılar}
% =====================================================================

\begin{lstlisting}
#include <stdio.h>
#include <time.h>
#include <unistd.h>

int main(void) {
    time_t simdi = time(NULL);               // 1970'ten beri saniye
    printf("Unix zaman: %ld\n", (long)simdi);
    printf("Yerel saat: %s", ctime(&simdi)); // insan-okur bicim

    struct tm *t = localtime(&simdi);        // parcalara ayir
    printf("Yil: %d, Ay: %d, Gun: %d\n",
           t->tm_year + 1900, t->tm_mon + 1, t->tm_mday);

    // Yuksek cozunurluklu sure olcumu (monotonik saat)
    struct timespec bas, son;
    clock_gettime(CLOCK_MONOTONIC, &bas);
    usleep(150000);                          // 150 ms uyu (mikrosaniye)
    clock_gettime(CLOCK_MONOTONIC, &son);

    double gecen = (son.tv_sec - bas.tv_sec)
                 + (son.tv_nsec - bas.tv_nsec) / 1e9;
    printf("Gecen sure: %.3f sn\n", gecen);
    return 0;
}
\end{lstlisting}

\ipucu{Süre ölçümünde \code{CLOCK\_MONOTONIC} kullan; \code{CLOCK\_REALTIME}
(duvar saati) NTP ya da kullanıcı tarafından geri alınabilir ve negatif
süreler ölçülmesine yol açabilir.}

% =====================================================================
\gsec{GDB ile Hata Ayıklama}
% =====================================================================

GDB (GNU Debugger), çalışan bir programı durdurup adım adım yürütmeni,
değişkenleri incelemeni, çökme yerini bulmanı ve belleği okumanı sağlar.
Segfault, bellek bozulması ve mantık hatalarını bulmanın en güçlü aracıdır.

\gsub{Hazırlık: Doğru Derleme}

\begin{lstlisting}[style=shstyle]
# -g: hata ayiklama sembolleri (degisken adlari, satir no)
# -O0: optimizasyonu kapat (kod satirlarla ortusur, degisken kaybolmaz)
$ gcc -g -O0 -Wall program.c -o program

# Cekirdek (core) dump'lari etkinlestir (cokme sonrasi analiz icin)
$ ulimit -c unlimited

# GDB'yi baslat
$ gdb ./program              # programla baslat
$ gdb --args ./program a b   # argumanlarla baslat
$ gdb -tui ./program         # TUI (metin arayuzu) ile baslat
\end{lstlisting}

\gsub{Programı Çalıştırma ve Durdurma}

\begin{center}
\begin{tabularx}{\textwidth}{@{}>{\ttfamily\color{anaMavi}}l l >{\raggedright\arraybackslash}X@{}}
\toprule
\rowcolor{acikMavi}\color{anaMavi}\textbf{Komut} & \color{anaMavi}\textbf{Kısa} & \color{anaMavi}\textbf{Görev}\\
\midrule
run [arg] & r & Programı baştan çalıştır.\\
\rowcolor{acikGri} start & --- & main'de durup çalıştır (geçici kesme).\\
continue & c & Sonraki kesmeye kadar devam et.\\
\rowcolor{acikGri} kill & k & Çalışan programı durdur.\\
quit & q & GDB'den çık.\\
\rowcolor{acikGri} Ctrl+C & --- & Çalışan programı anında durdur.\\
\bottomrule
\end{tabularx}
\end{center}

\gsub{Kesme Noktaları (Breakpoints)}

\begin{center}
\begin{tabularx}{\textwidth}{@{}>{\ttfamily\color{anaMavi}}l >{\raggedright\arraybackslash}X@{}}
\toprule
\rowcolor{acikMavi}\color{anaMavi}\textbf{Komut} & \color{anaMavi}\textbf{Görev}\\
\midrule
break main & \code{main} fonksiyonuna kesme koy (\code{b main}).\\
\rowcolor{acikGri} break dosya.c:42 & 42. satıra kesme koy.\\
break fonk if x>5 & \emph{Koşullu} kesme: yalnızca x>5 iken dur.\\
\rowcolor{acikGri} tbreak & Tek seferlik (geçici) kesme.\\
watch degisken & \emph{Watchpoint}: değer her değiştiğinde dur.\\
\rowcolor{acikGri} rwatch / awatch & Okunduğunda / okunup-yazıldığında dur.\\
info breakpoints & Tüm kesmeleri listele (\code{i b}).\\
\rowcolor{acikGri} delete 2 / disable 2 & 2 no'lu kesmeyi sil / devre dışı bırak.\\
\bottomrule
\end{tabularx}
\end{center}

\bilgi{\textbf{Watchpoint}, bellek bozulması hatalarını bulmanın en güçlü
yoludur: "bu değişken beklenmedik biçimde değişti" derken \code{watch degisken}
koy, GDB değeri değiştiren \emph{tam satırda} durur.}

\gsub{Adım Adım Yürütme (Stepping)}

\begin{center}
\begin{tabularx}{\textwidth}{@{}>{\ttfamily\color{anaMavi}}l l >{\raggedright\arraybackslash}X@{}}
\toprule
\rowcolor{acikMavi}\color{anaMavi}\textbf{Komut} & \color{anaMavi}\textbf{Kısa} & \color{anaMavi}\textbf{Görev}\\
\midrule
next & n & Sonraki satır (fonksiyon çağrısının \emph{üstünden} atla).\\
\rowcolor{acikGri} step & s & Sonraki satır (fonksiyonun \emph{içine} gir).\\
finish & fin & Mevcut fonksiyondan çık (dönüş değerini göster).\\
\rowcolor{acikGri} until 50 & u 50 & 50. satıra kadar çalıştır (döngüden çıkmak için).\\
nexti / stepi & ni/si & Tek makine komutu ilerle (assembly seviyesi).\\
\rowcolor{acikGri} return & --- & Fonksiyondan hemen dön (geri kalanı atla).\\
\bottomrule
\end{tabularx}
\end{center}

\uyari{\code{next} ile \code{step} farkı kritiktir: bir fonksiyon çağrısı
satırında \code{next} onu tek adımda geçer, \code{step} ise içine dalar.
Kütüphane fonksiyonlarına dalmamak için genelde \code{next} kullanılır.}

\gsub{Veriyi İnceleme (Print, Examine)}

\begin{center}
\begin{tabularx}{\textwidth}{@{}>{\ttfamily\color{anaMavi}}l >{\raggedright\arraybackslash}X@{}}
\toprule
\rowcolor{acikMavi}\color{anaMavi}\textbf{Komut} & \color{anaMavi}\textbf{Görev}\\
\midrule
print x & \code{x}'in değerini yazdır (\code{p x}).\\
\rowcolor{acikGri} print/x x & Onaltılık (hex) yazdır (\code{/d} ondalık, \code{/c} karakter, \code{/t} ikilik).\\
print dizi@10 & Dizinin ilk 10 elemanını yazdır.\\
\rowcolor{acikGri} print *ptr & İşaretçinin gösterdiği değeri yazdır.\\
print yapi & Bir struct'ın tüm alanlarını yazdır.\\
\rowcolor{acikGri} display x & Her durakta \code{x}'i otomatik göster.\\
x/8xw \&dizi & Belleği incele: 8 word, hex (\code{x/} = examine).\\
\rowcolor{acikGri} set var x = 5 & Değişkeni çalışırken değiştir.\\
ptype x / whatis x & Tipini göster.\\
\bottomrule
\end{tabularx}
\end{center}

\notk[x/ (examine) biçimi]{\code{x/NFU adres} --- \textbf{N}=adet,
\textbf{F}=biçim (x=hex, d=ondalık, c=karakter, s=string, i=komut),
\textbf{U}=birim (b=bayt, h=2, w=4, g=8 bayt). Örnek: \code{x/16xb ptr} =
16 baytı hex olarak; \code{x/s ptr} = string olarak.}

\gsub{Çağrı Yığını (Backtrace)}

Segfault olduğunda ilk yapılacak: \code{backtrace} ile "nereden nasıl
buraya gelindi" zincirini görmek.

\begin{center}
\begin{tabularx}{\textwidth}{@{}>{\ttfamily\color{anaMavi}}l l >{\raggedright\arraybackslash}X@{}}
\toprule
\rowcolor{acikMavi}\color{anaMavi}\textbf{Komut} & \color{anaMavi}\textbf{Kısa} & \color{anaMavi}\textbf{Görev}\\
\midrule
backtrace & bt & Tüm çağrı yığınını göster.\\
\rowcolor{acikGri} bt full & --- & Yığın + her çerçevedeki yerel değişkenler.\\
frame 2 & f 2 & 2 no'lu çerçeveye geç.\\
\rowcolor{acikGri} up / down & --- & Bir üst / alt çerçeveye geç.\\
info locals & i lo & Yereldeki tüm değişkenleri göster.\\
\rowcolor{acikGri} info args & i ar & Fonksiyon argümanlarını göster.\\
\bottomrule
\end{tabularx}
\end{center}

\ornek[Segfault'u bulma oturumu]{Klasik akış: çalıştır, çök, backtrace,
suçlu satırı ve değişkeni incele.}
\begin{lstlisting}[style=shstyle]
(gdb) run
Program received signal SIGSEGV, Segmentation fault.
0x0000... in listeye_ekle (bas=0x0, veri=5) at liste.c:14
14          bas->sonraki = yeni;
(gdb) bt                      # cagri yigini
#0  listeye_ekle (bas=0x0, ...) at liste.c:14
#1  0x... in main () at liste.c:30
(gdb) print bas               # suclu: bas NULL!
$1 = (Dugum *) 0x0
(gdb) frame 1                 # cagirana cik
(gdb) print liste             # neden NULL gecilmis?
\end{lstlisting}

\gsub{TUI: Metin Kullanıcı Arayüzü ve Layout'lar}

TUI modu, kaynağı/assembly'yi/register'ları GDB komut satırının üstünde
bir pencerede gösterir --- adım adım yürütmeyi görsel yapar.

\begin{center}
\begin{tabularx}{\textwidth}{@{}>{\ttfamily\color{anaMavi}}l >{\raggedright\arraybackslash}X@{}}
\toprule
\rowcolor{acikMavi}\color{anaMavi}\textbf{Komut / Kısayol} & \color{anaMavi}\textbf{Görev}\\
\midrule
Ctrl+X sonra A & TUI modunu aç/kapat.\\
\rowcolor{acikGri} layout src & Kaynak kodu penceresi.\\
layout asm & Assembly penceresi.\\
\rowcolor{acikGri} layout split & Kaynak + assembly birlikte.\\
layout regs & Register penceresi (+ kaynak).\\
\rowcolor{acikGri} Ctrl+X sonra 2 & Pencere düzenini değiştir.\\
Ctrl+L & Ekranı yenile (bozulursa).\\
\rowcolor{acikGri} Ctrl+P / Ctrl+N & Komut geçmişinde geri/ileri.\\
focus cmd / focus src & Klavye odağını komut/kaynak penceresine ver.\\
\bottomrule
\end{tabularx}
\end{center}

\ipucu{TUI'de yön tuşları odaklanan pencereyi kaydırır. Komut satırında geçmişe
erişmek için \code{focus cmd} ile odağı komut penceresine al; kaynağı kaydırmak
için \code{focus src}. Ekran bozulursa \code{Ctrl+L} yeniler.}

\gsub{Core Dump ve Çalışan Sürece Bağlanma}

\begin{lstlisting}[style=shstyle]
# 1) Cokme sonrasi (post-mortem) analiz: core dosyasiyla
$ ulimit -c unlimited            # once etkinlestir
$ ./program                      # cokup "core" uretir
$ gdb ./program core             # core'u yukle
(gdb) bt                         # coktugu andaki yigin

# 2) Zaten calisan bir surece baglanma (PID ile)
$ gdb -p 12345                   # 12345 PID'li surece baglan
(gdb) bt                         # o anki durumu incele
(gdb) detach                     # birak, surec calismaya devam etsin
\end{lstlisting}

\gsub{Diğer Faydalı Komutlar}

\begin{center}
\begin{tabularx}{\textwidth}{@{}>{\ttfamily\color{anaMavi}}l >{\raggedright\arraybackslash}X@{}}
\toprule
\rowcolor{acikMavi}\color{anaMavi}\textbf{Komut} & \color{anaMavi}\textbf{Görev}\\
\midrule
list / list 40 & Kaynak kodu göster (\code{l}).\\
\rowcolor{acikGri} info registers & Tüm register değerleri (\code{i r}).\\
info functions / info variables & Sembolleri listele.\\
\rowcolor{acikGri} call fonk(3) & Program içindeki bir fonksiyonu elle çağır.\\
break + commands & Kesmeye ulaşınca otomatik komut dizisi çalıştır.\\
\rowcolor{acikGri} set print pretty on & Struct'ları girintili/okunur yazdır.\\
tbreak main; run & main'de dur (start ile aynı).\\
\bottomrule
\end{tabularx}
\end{center}

\ornek[İş parçacığı hata ayıklama]{Çok iş parçacıklı programlarda thread'ler
arası geçiş.}
\begin{lstlisting}[style=shstyle]
(gdb) info threads          # tum iplikleri listele
(gdb) thread 3              # 3 no'lu iplige gec
(gdb) thread apply all bt   # TUM ipliklerin yiginini goster (deadlock avi)
\end{lstlisting}

\gsub{.gdbinit ile Otomasyon}

Ev dizinindeki \code{\textasciitilde/.gdbinit} dosyası her oturumda çalışır;
sık kullanılan ayarları ve makroları kalıcı yapar.

\begin{lstlisting}[style=shstyle]
# ~/.gdbinit ornegi
set print pretty on          # struct'lari okunur yazdir
set pagination off           # "--More--" duraklamasini kapat
set history save on          # komut gecmisini sakla

# Ozel komut (kisayol) tanimla:
define bt_full
    thread apply all backtrace full
end
\end{lstlisting}

\bilgi{GDB alternatifleri: daha modern bir arayüz için \textbf{gdb -tui} ya da
\textbf{cgdb}, Python tabanlı zengin görünüm için \textbf{gdb-dashboard} veya
\textbf{pwndbg}/\textbf{GEF} eklentileri kullanılabilir. Bellek hataları için
GDB'yi \textbf{AddressSanitizer} (\code{-fsanitize=address}) ile birlikte
kullanmak, hatayı oluştuğu anda yakalar.}

% =====================================================================
\gsec{Ek: Derleme, Hata Ayıklama ve Araçlar}
% =====================================================================

\begin{center}
\begin{tabularx}{\textwidth}{@{}>{\ttfamily\color{anaMavi}}l >{\raggedright\arraybackslash}X@{}}
\toprule
\rowcolor{acikMavi}\color{anaMavi}\textbf{Komut} & \color{anaMavi}\textbf{Görev}\\
\midrule
gcc -Wall -Wextra -g & Uyarılar + hata ayıklama sembolleri.\\
\rowcolor{acikGri} gcc -fsanitize=address & Bellek hatalarını çalışma anında yakala.\\
gcc -pthread & POSIX threads bağla.\\
\rowcolor{acikGri} gdb ./program & Etkileşimli hata ayıklayıcı.\\
valgrind ./program & Bellek sızıntısı/hatası dedektörü.\\
\rowcolor{acikGri} strace ./program & Yapılan sistem çağrılarını izle.\\
ltrace ./program & Kütüphane çağrılarını izle.\\
\rowcolor{acikGri} man 2 / man 3 & Sistem çağrısı / kütüphane el kitabı.\\
\bottomrule
\end{tabularx}
\end{center}

\ornek[Tipik bir hata ayıklama oturumu]{Segfault'u \code{gdb} ile bulmak.}
\begin{lstlisting}[style=shstyle]
$ gcc -g -O0 program.c -o program        # -g sart, -O0 optimizasyonu kapat
$ gdb ./program
(gdb) run                                # cokene kadar calistir
(gdb) backtrace                          # cagri yigini: nerede coktu?
(gdb) frame 1                            # bir ust cerceveye gec
(gdb) print degisken                     # bir degiskenin degerini gor
(gdb) break dosya.c:42                   # 42. satira kesme noktasi
(gdb) continue                           # devam et
\end{lstlisting}

\bilgi{\textbf{Öğrenmeye devam:} \code{man} sayfaları en güvenilir kaynaktır.
Bu rehberdeki her sistem çağrısının detaylı davranışı, hata kodları ve
örnekleri için \code{man 2 <ad>} komutuna başvur. Kod örneklerini kendi
makinende derleyip \code{strace} ile hangi sistem çağrılarını yaptıklarını
izlemek, sistem programlamayı gerçekten kavramanın en iyi yoludur.}

\vfill
\begin{center}
{\color{ortaGri}\rule{0.5\linewidth}{0.6pt}}\\[6pt]
{\color{ortaGri}\small Rehberin sonu --- İyi kodlamalar!}
\end{center}

\end{document}
